\documentclass[11pt]{article}
\usepackage[T1]{fontenc}
\usepackage{lmodern,amsmath,amssymb,amsthm,mathtools}
\usepackage[margin=1in]{geometry}
\usepackage{microtype,enumitem,booktabs,xcolor}
\usepackage[numbers,sort&compress]{natbib}
\usepackage{hyperref}
\hypersetup{colorlinks=true,linkcolor=blue!45!black,citecolor=blue!45!black,
urlcolor=blue!45!black,pdftitle={On the structural conjecture for extremizers of the independent Berry-Esseen inequality},pdfauthor={Logan Bell}}
\setlength{\emergencystretch}{2em}
\numberwithin{equation}{section}
\newtheorem{theorem}{Theorem}[section]
\newtheorem{proposition}[theorem]{Proposition}
\newtheorem{lemma}[theorem]{Lemma}
\newtheorem{corollary}[theorem]{Corollary}
\theoremstyle{definition}
\newtheorem{definition}[theorem]{Definition}
\newtheorem{conjecture}[theorem]{Conjecture}
\theoremstyle{remark}
\newtheorem{remark}[theorem]{Remark}
\DeclareMathOperator{\Var}{Var}
\DeclareMathOperator{\supp}{supp}
\DeclareMathOperator{\Bern}{Bern}
\DeclareMathOperator{\Bin}{Bin}
\DeclareMathOperator{\Pois}{Pois}
\DeclareMathOperator{\Law}{Law}
\DeclareMathOperator{\dist}{dist}
\newcommand{\E}{\mathbb E}
\newcommand{\Pp}{\mathbb P}
\newcommand{\R}{\mathbb R}
\newcommand{\Z}{\mathbb Z}
\newcommand{\N}{\mathbb N}
\newcommand{\ind}{\mathbf1}
\newcommand{\ce}{C_{\mathrm E}}
\newcommand{\Cind}{C_{\mathrm{ind}}}
\newcommand{\pe}{p_{\mathrm E}}
\newcommand{\qe}{q_{\mathrm E}}
\newcommand{\phz}{\phi(0)}
\newcommand{\eps}{\varepsilon}
\newcommand{\ii}{\mathrm i}

\title{On the structural conjecture for extremizers\\of the independent Berry--Esseen inequality}
\author{Logan Bell}
\date{18 October 2026}

\begin{document}
\maketitle

\begin{abstract}
In the proved regimes of the independent Berry--Esseen problem---including
single summands, common-diameter two-point arrays, pairs of two-point
summands with arbitrary diameters, and general arrays with sufficiently small
Lyapunov ratio---arrays of independent centered random variables that maximize
the Berry--Esseen ratio consist of $m\le n$ identically distributed two-point
random variables together with $n-m$ deterministic zero variables. In this
paper we
formulate the structural and value conjectures for fixed numbers of summands
$n$, separate the determination of the optimal constant from the
classification and uniqueness of maximizing arrays, and examine what current
theorems and computations establish. First, we separate the choice of the
marginal distribution $\mu$ of the sum $S=X_1+\cdots+X_n$ from the choice of
independent summands $(X_1,\ldots,X_n)$ that realize $\mu$. Studying the
minimum third-moment functional $\mathcal M_n(\mu)=\inf\sum_{i=1}^n\E|X_i|^3$
over $n$-summand independent factorizations of $\mu$, we prove that
$\mathcal M_n(\mu)$ is attained on a $W_3^n$-compact fiber and is lower
semicontinuous, but for $n\ge2$ is continuous at $\mu$ if and only if $\mu$
is convolution-indecomposable. Using an elementary third-moment splitting
identity and probability generating functions, we solve the factorization
problem for Poisson-binomial and binomial sum distributions and deduce that
an $n$-summand global maximizer consists of $m\le n$ identically distributed
two-point summands and $n-m$ zeros if and only if its sum $S$ has a centered
binomial distribution of order $m\le n$. Second, we establish a sharp
two-parameter moment inequality on separated supports and an endogenous
characteristic-function criterion whose zero set over all finite or countable
independent arrays consists precisely of finite identically distributed
two-point arrays with zero padding. Third, we synthesize the partial
classifications proved in the companion manuscripts, explain what the stored
high-precision binomial scans and the certified four-summand unequal-diameter
local maximizer demonstrate, and isolate the missing global implication and
analytical obstructions that leave the general finite-$n$ conjectures open.
\end{abstract}

\tableofcontents

\section{Introduction: three questions on independent summands}\label{sec:intro}

\subsection{The Berry--Esseen ratio and the padded binomial family}

Let $X_1,\ldots,X_n$ be independent centered real random variables with
positive total variance and finite third absolute moments. Write
\begin{equation}\label{eq:basic-quantities}
 S=\sum_{i=1}^n X_i,\qquad
 V=\sum_{i=1}^n\E X_i^2>0,\qquad
 R=\sum_{i=1}^n\E|X_i|^3<\infty,\qquad
 L=\frac{R}{V^{3/2}},
\end{equation}
where $L$ is the Lyapunov ratio. Deterministic zero summands ($X_i=0$ almost
surely) are permitted; they contribute zero to $S$, $V$, and $R$, and merely
leave available coordinate slots unused. Let
\begin{equation}\label{eq:kolmogorov-ratio}
 \Delta(S)=\sup_{x\in\R}\left|\Pp\left\{\frac{S}{\sqrt V}\le x\right\}-\Phi(x)\right|,\qquad
 \mathcal C(X_1,\ldots,X_n)=\frac{\Delta(S)}{L}=\frac{V^{3/2}\Delta(S)}{R},
\end{equation}
where $\Phi(x)=\int_{-\infty}^x\phi(u)\,du$ and
$\phi(u)=(2\pi)^{-1/2}e^{-u^2/2}$ are the standard normal distribution
function and density. Notice that $\mathcal C(X_1,\ldots,X_n)$ is invariant
under multiplying all summands by a common nonzero constant, under common
sign reflection $X_i\mapsto -X_i$, and under permuting coordinates.

For each integer $n\ge1$, let
\begin{equation}\label{eq:cn-cind}
 C_n=\sup\mathcal C(X_1,\ldots,X_n),\qquad
 \Cind=\sup_{n\ge1}C_n,
\end{equation}
where the supremum in $C_n$ is taken over all $n$-tuples of independent
centered real random variables satisfying \eqref{eq:basic-quantities}.
Because an $n$-summand array can be extended to $n+1$ summands by adjoining
$X_{n+1}=0$, the sequence $(C_n)_{n\ge1}$ is nondecreasing. As we recall in
Section~\ref{sec:known}, the supremum $C_n$ is attained for every $n\ge1$,
and requiring every summand to have strictly positive variance yields the
exact same supremum $C_n$ (though the positive-variance supremum need not be
attained, as occurs at $n=2$ when the summands are restricted to two-point
distributions \citep[Theorem~1.2]{BellBernoulli2026}).

A fundamental special family consists of arrays in which $m\in\{1,\ldots,n\}$
summands are identically distributed centered two-point (Bernoulli) random
variables and the remaining $n-m$ summands are zero:
\begin{equation}\label{eq:padded-binomial-array}
 X_i=h(B_i-p)\quad(1\le i\le m),\qquad
 X_i=0\quad(m<i\le n),
\end{equation}
where $h>0$ is the common support diameter, $B_1,\ldots,B_m$ are independent
$\Bern(p)$ random variables with $\Pp(B_i=1)=p\in(0,1)$ and
$\Pp(B_i=0)=q:=1-p$. For the array \eqref{eq:padded-binomial-array},
\[
 V=mh^2pq,\qquad
 R=mh^3pq(p^2+q^2),\qquad
 L=\frac{p^2+q^2}{\sqrt{mpq}}.
\]
Writing $K=\sum_{i=1}^m B_i\sim\Bin(m,p)$ and $z_k=(k-mp)/\sqrt{mpq}$, the
Kolmogorov distance $\Delta(S)$ is achieved at one of the jump points
$k\in\{0,1,\ldots,m\}$ on either the right-continuous or the left-continuous
branch of the distribution function, so the Berry--Esseen ratio of
\eqref{eq:padded-binomial-array} is
\begin{equation}\label{eq:rmp}
 R_m(p)=\frac{\sqrt{mpq}}{p^2+q^2}
 \max_{0\le k\le m}\bigl\{\Pp(K\le k)-\Phi(z_k),\,\Phi(z_k)-\Pp(K<k)\bigr\}.
\end{equation}
Since $R_m(p)=R_m(1-p)$ grows like $\sqrt{mp}$ as $p\downarrow0$, the
continuous function $p\mapsto R_m(p)$ attains its maximum on $(0,1)$. Define
the $m$-summand binomial maximum and the $n$-summand padded binomial envelope
by
\begin{equation}\label{eq:bm-pn}
 B_m=\max_{0<p<1}R_m(p),\qquad
 P_n=\max_{1\le m\le n}B_m.
\end{equation}
Clearly $P_n\le C_n$ for every $n\ge1$.

In 1956, \citet{Esseen1956} proved that among identically distributed
summands with finite third absolute moments, the largest asymptotic value of
$\Delta/L$ as $m\to\infty$ is
\begin{equation}\label{eq:esseen-constants}
 \ce=\frac{3+\sqrt{10}}{6\sqrt{2\pi}}=0.40973218370239634299\ldots,
\end{equation}
and that this limit is attained by two-point summands with success
probability
\begin{equation}\label{eq:esseen-probability}
 \pe=\frac{4-\sqrt{10}}{2}=0.41886116991581033400\ldots
\end{equation}
or $\qe=1-\pe=( \sqrt{10}-2)/2$. Sixty years later, \citet[Theorem~1]{Schulz2016}
proved that $B_m<\ce$ for every $m\ge1$, so that $\sup_{n\ge1}P_n=\lim_{m\to\infty}B_m=\ce$.

\subsection{Separating value, structure, and uniqueness}

The central theme of this paper is the empirical and theoretical observation
that arrays maximizing $\mathcal C(X_1,\ldots,X_n)$ appear always to belong
to the padded binomial family \eqref{eq:padded-binomial-array}. To analyze
this phenomenon precisely, one must keep three questions distinct:

\begin{enumerate}[label=\textbf{Question \arabic*.},ref=\arabic*,leftmargin=6.5em]
\item\label{q:value} \textbf{(Optimal value)}
\emph{Does $C_n=P_n$ hold for a given $n\ge1$, and does $\Cind=\ce$ hold
across all $n$?}

This is an inequality statement: it asks whether $\Delta(S)\le P_n L$ holds
for every $n$-summand array (and $\Delta(S)\le\ce L$ for all finite arrays).
Because $P_n<\ce$ for every $n$ and $\lim_{n\to\infty}P_n=\ce$, the equality
$C_n=P_n$ for all $n\ge1$ would immediately imply $\Cind=\ce$, whereas
$\Cind=\ce$ is weaker than $C_n=P_n$ at finite $n$.

\item\label{q:structure} \textbf{(Classification of maximizing arrays)}
\emph{Does every $n$-summand array attaining $C_n$ have the form
\eqref{eq:padded-binomial-array} up to common scaling, reflection, and
coordinate permutation?}

Proving an upper bound on $C_n$, or even proving $C_n=P_n$, does not by
itself classify all arrays that attain the maximum; conversely, knowing that
every maximizer of $C_n$ has the form \eqref{eq:padded-binomial-array}
immediately implies $C_n=P_n$. Classification requires showing both that
every nonzero summand is supported on two points and that all nonzero
summands share the same support diameter $h$ and the same success probability
$p$ (with the same orientation).

\item\label{q:uniqueness} \textbf{(Uniqueness of the binomial parameters)}
\emph{Within the padded binomial family \eqref{eq:padded-binomial-array}, is
the maximizing active count $m_n\in\{1,\ldots,n\}$ in
$P_n=\max_{1\le m\le n}B_m$ unique, and for a given $m$, is the maximizing
probability $p_m\in(0,1/2]$ in $B_m=\max_{0<p\le1/2}R_m(p)$ unique?}

Even after an $n$-summand maximizer is shown to belong to the family
\eqref{eq:padded-binomial-array}, determining how many summands are zero and
whether two distinct probabilities or two distinct active counts $m\ne m'$
can tie for the maximum $P_n$ is a separate analytical problem in the
one-variable family $R_m(p)$.
\end{enumerate}

Why can a maximizer at horizon $n$ set $n-m>0$ coordinates to zero instead of
using all $n$ available summands? As already occurs at $n=2$, where
$B_2\approx0.355891<B_1\approx0.370352$, and as persists for arbitrarily
large $n$, the Kolmogorov distance of a binomial sum $\Bin(m,p)$ is evaluated
at an integer threshold $k\in\{0,\ldots,m\}$. When $m$ is large, the
optimal probability $p$ lies within $O(m^{-1})$ of Esseen's probability
$\pe$, and the optimal fractional offset $mp-k$ between the mean and the
integer threshold converges to $c_*=(20+\sqrt{10})/72\approx0.321698$.
Because $k$ must be an integer, $mp-k$ cannot be well aligned with $c_*$ near
$p=\pe$ at every integer $m$. Since $B_m=\ce-Q_*/m+O(m^{-2})$, reducing the
active sample size from $n$ to $m=n-j$ ($j\in\{0,1,2,3,4\}$) incurs an
$O(n^{-2})$ penalty $Q_*j/n^2$ in the leading $-Q_*/m$ term while shifting
the fractional lattice phase $(n-j)\pe\bmod1$ to reduce the $O(n^{-2})$
threshold-quantization loss, strictly increasing the Berry--Esseen ratio
whenever $n\pe\bmod1$ lies in specific phase intervals
\citep[Section~4]{BellBernoulli2026}.

\subsection{Outline of the paper}

The purpose of this paper is to give a self-contained account of the
structural theory surrounding Questions~\ref{q:value}--\ref{q:uniqueness},
connecting the theorems proved in the seven companion manuscripts
\citep{BellIndependent2026,BellFinite2026,BellSmall2026,BellTransport2026,BellBernoulli2026,BellSmallStructure2026,BellExtremizers2026}
with two exact structural mechanisms and with the repository's computational
records:

\begin{itemize}[leftmargin=2em]
\item In \textbf{Section~\ref{sec:factorization}} we examine how choosing the
marginal distribution $\mu=\Law(S)$ of the sum differs from choosing the
independent summands $(X_1,\ldots,X_n)$ that produce $\mu$. The numerator
$\Delta(S)$ and total variance $V$ depend only on $\mu$, so for any non-normal
sum distribution $\mu$ (where $\Delta(\mu)>0$), maximizing
$\mathcal C(X_1,\ldots,X_n)$ among independent factorizations of $\mu$ is
equivalent to minimizing the third-moment sum $\sum_{i=1}^n\E|X_i|^3$ over all
$n$-summand independent factorizations of $\mu$. Using an elementary
third-moment splitting identity (Proposition~\ref{prop:splitting}), we prove
that the minimum third-moment functional $\mathcal M_n(\mu)$ is attained and
lower semicontinuous in $W_3$, but for $n\ge2$ is continuous at $\mu$ if and
only if $\mu$ is convolution-indecomposable (Theorem~\ref{thm:fiber-regularity}).
Using probability generating functions, we solve the factorization problem
explicitly when $\mu$ is a Poisson-binomial or binomial law and deduce that
an $n$-summand maximizer of $C_n$ has the padded i.i.d.\ two-point form
\eqref{eq:padded-binomial-array} if and only if its sum $S$ has a centered
binomial distribution of order $m\le n$ (Theorem~\ref{thm:binomial-fiber}).

\item In \textbf{Section~\ref{sec:criteria}} we establish exact moment and
characteristic-function criteria that single out padded i.i.d.\ two-point
arrays across all finite or countable independent arrays. First, for
summands whose supports have spacing at least $h>0$, a size-biased
decomposition yields a sharp two-parameter family of linear inequalities
between $R$, $hV$, and the signed third moment $K_3=\sum_i\E X_i^3$, with
equality if and only if all nonzero summands share a single two-point law
(Theorem~\ref{thm:separated-moment}). Second, without assuming discreteness or
support separation, the moments $V$, $R$, and $K_3$ define a canonical scale
$\widehat h=(R+D)/V$ and effective active count $\widehat m=2V^3/(D(R+D))\le n$
such that the characteristic-function deficit
$\mathcal J=1-|\E e^{2\pi\ii S/\widehat h}|^2$ vanishes if and only if the
array is a finite i.i.d.\ two-point core with zero padding
(Theorem~\ref{thm:fourier-criterion}). We then compare these criteria with
the zero-bias Wasserstein transport inequality $2VW_1(S,S^*)\le R$ of
\citet{BellTransport2026}.

\item In \textbf{Section~\ref{sec:known}} we review the quantitative bounds
and partial structural classifications proved in the companion papers,
showing how the general three-point variational reduction
\citep{BellExtremizers2026}, the common-diameter Bernoulli classification and
phase expansion \citep{BellBernoulli2026}, and the small-Lyapunov global
maximizer classification \citep{BellSmallStructure2026} fit together.

\item In \textbf{Section~\ref{sec:computations}} we explain what the stored
computational datasets actually establish, distinguishing the floating-point
binomial evaluations through $m=3000$ and compound-Poisson comparisons from
the rigorous 320-bit interval certificate for a strict local maximizer at
$n=4$ with unequal support diameters $h=(1,1,1,2)$.

\item In \textbf{Section~\ref{sec:open}} we formulate the open finite-$n$
value, classification, and uniqueness conjectures and discuss the specific
analytical obstacles that remain between the proved regimes and the general
problem.
\end{itemize}

\section{Sum distributions versus independent summands}\label{sec:factorization}

\subsection{The minimum third-moment functional on a convolution fiber}

In the Berry--Esseen ratio
\[
 \mathcal C(X_1,\ldots,X_n)=\frac{V^{3/2}\Delta(S)}{R},
\]
the total variance $V=\Var(S)$ and the Kolmogorov discrepancy $\Delta(S)$
depend only on the marginal distribution $\mu=\Law(S)$ of the sum. The
individual laws $\mu_i=\Law(X_i)$ affect $\mathcal C(X_1,\ldots,X_n)$ solely
through the denominator $R=\sum_{i=1}^n\E|X_i|^3$ and the constraint that
their convolution $\mu_1*\cdots*\mu_n$ equals $\mu$.

Let $\mathcal P_3^0(\R)$ denote the space of centered Borel probability
measures on $\R$ with finite third absolute moment, equipped with the
$L^3$-Wasserstein distance
\[
 W_3(\mu,\nu)=\inf\bigl\{\bigl(\E|U-W|^3\bigr)^{1/3}:\Law(U)=\mu,\ \Law(W)=\nu\bigr\}.
\]
For $\mu\in\mathcal P_3^0(\R)$, write
\[
 V(\mu)=\int_\R x^2\,\mu(dx),\qquad
 \rho(\mu)=\int_\R |x|^3\,\mu(dx),
\]
and when $V(\mu)>0$, write
\[
 \Delta(\mu)=\sup_{x\in\R}\left|\mu\bigl((-\infty,x]\bigr)-\Phi\!\left(\frac{x}{\sqrt{V(\mu)}}\right)\right|.
\]
For each integer $n\ge1$ and each $\mu\in\mathcal P_3^0(\R)$, define the
\emph{convolution fiber}
\begin{equation}\label{eq:fiber-def}
 \mathcal F_n(\mu)=\bigl\{(\mu_1,\ldots,\mu_n)\in\mathcal P_3^0(\R)^n:
 \mu_1*\cdots*\mu_n=\mu\bigr\}
\end{equation}
and the \emph{minimum third-moment functional}
\begin{equation}\label{eq:mn-def}
 \mathcal M_n(\mu)=\inf_{(\mu_1,\ldots,\mu_n)\in\mathcal F_n(\mu)}
 \sum_{i=1}^n\rho(\mu_i).
\end{equation}
Because the Dirac mass $\delta_0$ belongs to $\mathcal P_3^0(\R)$ with
$\rho(\delta_0)=0$, the tuple $(\mu,\delta_0,\ldots,\delta_0)$ always belongs
to $\mathcal F_n(\mu)$, so $\mathcal F_n(\mu)$ is nonempty and
$\mathcal M_n(\mu)$ is nonincreasing in $n$:
\begin{equation}\label{eq:mn-bounds}
 \frac{V(\mu)^{3/2}}{\sqrt n}\le \mathcal M_n(\mu)\le \rho(\mu).
\end{equation}
Indeed, for any $(\mu_1,\ldots,\mu_n)\in\mathcal F_n(\mu)$ with variances
$v_i=V(\mu_i)$, Lyapunov's inequality $\rho(\mu_i)\ge v_i^{3/2}$ and H\"older's
inequality $\sum_{i=1}^n v_i\le n^{1/3}\bigl(\sum_{i=1}^n v_i^{3/2}\bigr)^{2/3}$
give the lower bound in \eqref{eq:mn-bounds}, with equality if and only if
$v_1=\cdots=v_n=V(\mu)/n$ and each $\mu_i$ is supported on
$\{\pm\sqrt{V(\mu)/n}\}$ with equal probabilities $1/2$.

Decomposing the supremum in \eqref{eq:cn-cind} first over the sum distribution
$\mu=\Law(S)$ and then over its convolution fiber $\mathcal F_n(\mu)$ gives
\begin{equation}\label{eq:two-stage-opt}
 C_n=\sup_{\substack{\mu\in\mathcal P_3^0(\R)\\ V(\mu)>0}}
 \frac{V(\mu)^{3/2}\Delta(\mu)}{\mathcal M_n(\mu)}.
\end{equation}
In particular, if $(X_1,\ldots,X_n)$ is any array attaining $C_n$ with sum
distribution $\mu=\Law(X_1+\cdots+X_n)$, then since $C_n\ge C_1>0$ implies
$\Delta(\mu)>0$, the summand distributions $\mu_i=\Law(X_i)$ must minimize
$\sum_{i=1}^n\rho(\mu_i)$ over $\mathcal F_n(\mu)$:
\begin{equation}\label{eq:maximizer-minimizes-price}
 \sum_{i=1}^n\E|X_i|^3=\mathcal M_n(\mu).
\end{equation}

\subsection{The third-moment splitting identity}

To understand which factorizations minimize $\sum_{i=1}^n\rho(\mu_i)$ on a
fixed fiber $\mathcal F_n(\mu)$, the basic tool is the strict superadditivity
of the third absolute moment under convolution of independent centered
variables \citep[Theorem~(2.6)]{CoxKemperman1983,Ruzankin2014}. The following
elementary identity makes the exact convolution excess completely explicit
without requiring symmetry or decomposition into two-point mixtures.

\begin{proposition}[Third-moment splitting identity]\label{prop:splitting}
Let $X$ and $Y$ be independent centered real random variables with finite
third absolute moments, and set $m=\min(|X|,|Y|)$ and $M=\max(|X|,|Y|)$. Then
\begin{equation}\label{eq:splitting-identity}
 \E|X+Y|^3-\E|X|^3-\E|Y|^3
 =2\,\E\bigl[\ind_{\{XY<0\}}m^2(3M-m)\bigr]\ge0.
\end{equation}
Moreover, equality holds in \eqref{eq:splitting-identity} if and only if at
least one of $X$ and $Y$ is zero almost surely.
\end{proposition}

\begin{proof}
We first verify that for all real numbers $x,y\in\R$,
\begin{equation}\label{eq:pointwise-cubic}
 |x+y|^3-|x|^3-|y|^3-3xy(|x|+|y|)
 =2\ind_{\{xy<0\}}\min(|x|,|y|)^2\bigl(3\max(|x|,|y|)-\min(|x|,|y|)\bigr).
\end{equation}
If $xy\ge0$, then $|x+y|=|x|+|y|$, so
$|x+y|^3-|x|^3-|y|^3=3|x|\,|y|(|x|+|y|)=3xy(|x|+|y|)$, and both sides of
\eqref{eq:pointwise-cubic} vanish. If $xy<0$, then by symmetry we may assume
$a=|x|\ge b=|y|>0$, so that $xy=-ab$ and $|x+y|=a-b$; the left-hand side of
\eqref{eq:pointwise-cubic} is then
\[
 (a-b)^3-a^3-b^3+3ab(a+b)=2b^2(3a-b),
\]
which coincides with the right-hand side of \eqref{eq:pointwise-cubic}.

Since $\E|X|^3<\infty$ and $\E|Y|^3<\infty$, every term in
\eqref{eq:pointwise-cubic} is integrable when evaluated at $(X,Y)$. Because
$X$ and $Y$ are independent and centered,
\[
 \E\bigl[XY|X|\bigr]=\E\bigl[X|X|\bigr]\,\E Y=0,\qquad
 \E\bigl[XY|Y|\bigr]=\E X\,\E\bigl[Y|Y|\bigr]=0.
\]
Taking expectations in \eqref{eq:pointwise-cubic} therefore yields
\eqref{eq:splitting-identity}.

On the event $\{XY<0\}$ we have $M\ge m>0$, so $m^2(3M-m)\ge2m^3>0$. If
neither $X$ nor $Y$ is zero almost surely, then since $\E X=\E Y=0$, each of
$X$ and $Y$ takes both strictly positive and strictly negative values with
positive probability; by independence, $\Pp(XY<0)=\Pp(X>0)\Pp(Y<0)+\Pp(X<0)\Pp(Y>0)>0$,
and therefore the right-hand side of \eqref{eq:splitting-identity} is
strictly positive.
\end{proof}

\subsection{Compactness, attainment, and discontinuity of the fiber minimum}

A measure $\mu\in\mathcal P_3^0(\R)$ is called \emph{convolution-indecomposable}
(or simply \emph{indecomposable} in $\mathcal P_3^0(\R)$) if every
factorization $\mu=\mu_1*\mu_2$ with $\mu_1,\mu_2\in\mathcal P_3^0(\R)$ has
$\mu_1=\delta_0$ or $\mu_2=\delta_0$. By Proposition~\ref{prop:splitting},
splitting a decomposable summand into two nonconstant independent centered
factors preserves the distribution of the sum $S$ while strictly decreasing
the sum of third absolute moments. This strict decrease has a striking
consequence for the regularity of $\mathcal M_n(\mu)$ as a function of the
sum law $\mu$: although $\mathcal M_n$ is attained on every fiber and is
lower semicontinuous, it is discontinuous at every decomposable distribution
whenever $n\ge2$.

\begin{theorem}[Compactness, attainment, and continuity of $\mathcal M_n$]\label{thm:fiber-regularity}
Fix an integer $n\ge1$.
\begin{enumerate}[label=(\roman*),leftmargin=2em]
\item\label{it:fiber-compact} For every $W_3$-compact subset
$\mathcal K\subset\mathcal P_3^0(\R)$, the union of convolution fibers
$\bigcup_{\mu\in\mathcal K}\mathcal F_n(\mu)$ is compact in
$\mathcal P_3^0(\R)^n$ under the product $W_3$ topology. In particular, for
every $\mu\in\mathcal P_3^0(\R)$, the fiber $\mathcal F_n(\mu)$ is
$W_3^n$-compact, the infimum $\mathcal M_n(\mu)$ in \eqref{eq:mn-def} is
attained, and $\mu\mapsto\mathcal M_n(\mu)$ is lower semicontinuous on
$(\mathcal P_3^0(\R),W_3)$.
\item\label{it:indecomp-criterion} If $n\ge2$, then for every
$\mu\in\mathcal P_3^0(\R)$,
\begin{equation}\label{eq:mn-equals-rho}
 \mathcal M_n(\mu)=\rho(\mu)
 \quad\Longleftrightarrow\quad
 \mu\text{ is convolution-indecomposable}.
\end{equation}
\item\label{it:dense-indecomp} Convolution-indecomposable finitely supported
measures are dense in $(\mathcal P_3^0(\R),W_3)$. Consequently, for fixed
$n\ge2$, the functional $\mathcal M_n$ is continuous at
$\mu\in\mathcal P_3^0(\R)$ if and only if $\mu$ is convolution-indecomposable.
\item\label{it:ratio-continuity} On the open domain
$\{\mu\in\mathcal P_3^0(\R):V(\mu)>0\}$, the optimized Berry--Esseen ratio
\[
 \mathcal Q_n(\mu)=\frac{V(\mu)^{3/2}\Delta(\mu)}{\mathcal M_n(\mu)}
\]
is upper semicontinuous in $W_3$. For fixed $n\ge2$, $\mathcal Q_n$ is
continuous at $\mu$ if and only if either $\mu$ is convolution-indecomposable
or $\mu$ is a centered normal distribution $N(0,V)$.
\end{enumerate}
\end{theorem}

\begin{proof}
\ref{it:fiber-compact} Let $(X_1,\ldots,X_n)$ be independent centered random
variables with $\Law(X_i)=\mu_i$ and $\Law(S)=\mu\in\mathcal K$, where
$S=X_1+\cdots+X_n$. Because the summands are independent and centered,
$\E[S\mid X_i]=X_i$ almost surely for each $i\in\{1,\ldots,n\}$. For every
$A\ge0$, the function $f_A(x)=(|x|-A)_+^3$ is convex on $\R$, and
$|x|^3\ind_{\{|x|>2A\}}\le8f_A(x)$ since $|x|>2A$ implies $|x|-A>|x|/2$.
Conditional Jensen's inequality therefore gives
\begin{equation}\label{eq:convex-tail}
 \int_{\{|x|>2A\}}|x|^3\,\mu_i(dx)
 =\E\bigl[|X_i|^3\ind_{\{|X_i|>2A\}}\bigr]
 \le 8\,\E f_A(X_i)\le 8\,\E f_A(S)
 =8\int_{\{|x|>A\}}(|x|-A)^3\,\mu(dx).
\end{equation}
Since $\mathcal K$ is $W_3$-compact, its third absolute moments are bounded
and uniformly integrable, so the right-hand side of \eqref{eq:convex-tail}
tends to $0$ as $A\to\infty$ uniformly over $\mu\in\mathcal K$. Hence each
marginal coordinate projection of $\bigcup_{\mu\in\mathcal K}\mathcal F_n(\mu)$
is tight with uniformly integrable third absolute moments, and is therefore
relatively compact in $(\mathcal P_3^0(\R),W_3)$. Because convergence in
$W_3$ preserves zero mean and convolution $(\mu_1,\ldots,\mu_n)\mapsto\mu_1*\cdots*\mu_n$
is continuous in $W_3$ (by coupling each coordinate independently and using
Minkowski's inequality), $\bigcup_{\mu\in\mathcal K}\mathcal F_n(\mu)$ is
closed and hence compact in $W_3^n$. Taking $\mathcal K=\{\mu\}$ shows that
$\mathcal F_n(\mu)$ is compact, and since $(\mu_1,\ldots,\mu_n)\mapsto\sum_{i=1}^n\rho(\mu_i)$
is continuous on $W_3^n$, the infimum $\mathcal M_n(\mu)$ is attained. If
$\mu^{(k)}\to\mu$ in $W_3$, applying compactness to the $W_3$-compact set
$\mathcal K=\{\mu\}\cup\{\mu^{(k)}:k\ge1\}$ and extracting a convergent
subsequence of attaining tuples $(\mu_1^{(k)},\ldots,\mu_n^{(k)})\in\mathcal F_n(\mu^{(k)})$
yields $\mathcal M_n(\mu)\le\liminf_{k\to\infty}\mathcal M_n(\mu^{(k)})$.

\ref{it:indecomp-criterion} Let $n\ge2$. If $\mu=\nu_1*\nu_2$ with
$\nu_1,\nu_2\in\mathcal P_3^0(\R)\setminus\{\delta_0\}$, then
$(\nu_1,\nu_2,\delta_0,\ldots,\delta_0)\in\mathcal F_n(\mu)$ has third-moment
sum $\rho(\nu_1)+\rho(\nu_2)<\rho(\mu)$ by Proposition~\ref{prop:splitting},
so $\mathcal M_n(\mu)<\rho(\mu)$. Conversely, if $(\mu_1,\ldots,\mu_n)\in\mathcal F_n(\mu)$
has at least two nondegenerate components, say $\mu_1\ne\delta_0$ and
$\mu_2\ne\delta_0$, then grouping $\nu_1=\mu_1$ and
$\nu_2=\mu_2*\cdots*\mu_n$ gives a nontrivial factorization $\mu=\nu_1*\nu_2$
(since $\Var(\nu_2)=\sum_{i=2}^n\Var(\mu_i)\ge\Var(\mu_2)>0$), so $\mu$ is
decomposable. Thus when $\mu$ is indecomposable, every tuple in
$\mathcal F_n(\mu)$ has at most one nonzero component, which must equal $\mu$,
and hence $\mathcal M_n(\mu)=\rho(\mu)$.

\ref{it:dense-indecomp} Suppose $\nu\in\mathcal P_3^0(\R)$ is supported on a
finite set $S=\{x_1<\cdots<x_s\}$ with $\nu(\{x_j\})>0$ for all $j$, and
suppose that all $\binom{s}{2}$ positive differences $x_r-x_j$ ($1\le j<r\le s$)
are distinct. We claim that $\nu$ is convolution-indecomposable. Indeed, if
$\nu=\nu_1*\nu_2$ with neither factor degenerate, then $\supp(\nu_1)$ contains
at least two points $a_0<a_1$ and $\supp(\nu_2)$ contains at least two points
$b_0<b_1$. Because all four sums $a_0+b_0$, $a_1+b_0$, $a_0+b_1$, $a_1+b_1$
belong to $\supp(\nu)=S$, either $a_1+b_0=a_0+b_1$ (in which case the three
distinct points $a_0+b_0<a_1+b_0<a_1+b_1$ of $S$ repeat the positive
difference $a_1-a_0=b_1-b_0$) or the four sums are four distinct points of $S$
among which the positive difference $a_1-a_0=(a_1+b_0)-(a_0+b_0)=(a_1+b_1)-(a_0+b_1)$
occurs twice, contradicting the assumption that all positive differences in
$S$ are distinct.

Any $\mu\in\mathcal P_3^0(\R)$ can be approximated arbitrarily closely in
$W_3$ by finitely supported centered measures, and an arbitrarily small
perturbation of the support points outside the finite union of hyperplanes
where two coordinates or two pairwise differences coincide (followed by
subtracting the perturbed mean, which shifts all support points by the same
constant and preserves all pairwise differences) produces a sequence
$\mu^{(k)}\to\mu$ in $W_3$ of finitely supported indecomposable measures. For
fixed $n\ge2$, if $\mu$ is indecomposable, then for any sequence $\nu^{(k)}\to\mu$
in $W_3$, lower semicontinuity and the upper bound
$\mathcal M_n(\nu^{(k)})\le\rho(\nu^{(k)})\to\rho(\mu)=\mathcal M_n(\mu)$
force $\mathcal M_n(\nu^{(k)})\to\mathcal M_n(\mu)$. Conversely, if $\mu$ is
decomposable, then $\mathcal M_n(\mu)<\rho(\mu)$ by \ref{it:indecomp-criterion},
whereas along a sequence of indecomposable measures $\mu^{(k)}\to\mu$ in $W_3$
we have $\mathcal M_n(\mu^{(k)})=\rho(\mu^{(k)})\to\rho(\mu)>\mathcal M_n(\mu)$,
so $\mathcal M_n$ is discontinuous at $\mu$.

\ref{it:ratio-continuity} For any two centered probability measures $\mu,\nu$
with finite first moments and any coupling $(U,W)$ with
$\E|U-W|=W_1(\mu,\nu)$, centering gives
$\E(W-U)_+=\E(U-W)_+=\frac12\E|U-W|=\frac12 W_1(\mu,\nu)$. By Markov's
inequality and the bound $\|\Phi'\|_\infty=\phz$, for every $x\in\R$ and
$\delta>0$,
\[
 \Pp(U\le x)-\Phi(x)
 \le \Pp(W\le x+\delta)-\Phi(x+\delta)+\phz\delta+\frac{\E(W-U)_+}{\delta}.
\]
Choosing $\delta=\sqrt{W_1(\mu,\nu)/(2\phz)}$ and repeating the argument for
the opposite sign and with $U$ and $W$ interchanged yields
\[
 \left|\sup_{x\in\R}\bigl|\Pp(U\le x)-\Phi(x)\bigr|
      -\sup_{x\in\R}\bigl|\Pp(W\le x)-\Phi(x)\bigr|\right|
 \le\sqrt{2\phz W_1(\mu,\nu)}.
\]
Since standardization $\mu\mapsto\Law(S/\sqrt{V(\mu)})$ is continuous in $W_3$
on $\{V(\mu)>0\}$, the numerator $V(\mu)^{3/2}\Delta(\mu)$ is continuous in
$W_3$. Because $\mathcal M_n(\mu)\ge V(\mu)^{3/2}/\sqrt n>0$ is lower
semicontinuous by \ref{it:fiber-compact}, the quotient $\mathcal Q_n(\mu)$ is
upper semicontinuous.

If $\mu$ is indecomposable, both numerator and denominator of $\mathcal Q_n$
are continuous at $\mu$. If $\mu=N(0,V)$, then $\Delta(\mu)=0$, and as
$\nu\to\mu$ in $W_3$ the bound $\mathcal M_n(\nu)\ge V(\nu)^{3/2}/\sqrt n$
gives $0\le\mathcal Q_n(\nu)\le\sqrt n\,\Delta(\nu)\to0=\mathcal Q_n(\mu)$,
so $\mathcal Q_n$ is continuous at $\mu$ even though $\mathcal M_n$ jumps from
$\rho(N(0,V))$ along indecomposable approximants to
$\mathcal M_n(N(0,V))=\rho(N(0,V))/\sqrt n$. Finally, if $\mu$ is decomposable
and not normal, then $\Delta(\mu)>0$ while $\mathcal M_n(\mu)<\rho(\mu)$;
along indecomposable approximants $\mu^{(k)}\to\mu$ in $W_3$,
\[
 \lim_{k\to\infty}\mathcal Q_n(\mu^{(k)})
 =\frac{V(\mu)^{3/2}\Delta(\mu)}{\rho(\mu)}
 <\frac{V(\mu)^{3/2}\Delta(\mu)}{\mathcal M_n(\mu)}=\mathcal Q_n(\mu),
\]
so $\mathcal Q_n$ is discontinuous at $\mu$.
\end{proof}

Theorem~\ref{thm:fiber-regularity} explains why the reformulation
\eqref{eq:two-stage-opt} over marginal sum laws $\mu$ cannot be treated as a
smooth variational problem on $\mathcal P_3^0(\R)$: the denominator
$\mathcal M_n(\mu)$ drops downward precisely on the thin locus of
distributions $\mu$ that admit nontrivial independent factorizations.

\subsection{Factorizations of Poisson-binomial and binomial sum distributions}

When the sum distribution $\mu$ is a centered Poisson-binomial or binomial
distribution, we can determine all elements of the convolution fiber
$\mathcal F_n(\mu)$ explicitly and identify the unique factorizations that
achieve $\mathcal M_n(\mu)$.

\begin{theorem}[Factorizations of Poisson-binomial and binomial laws]\label{thm:binomial-fiber}
Fix an integer $n\ge1$ and a lattice span $g>0$.
\begin{enumerate}[label=(\roman*),leftmargin=2em]
\item\label{it:pb-fiber} Let $\mu=\Law\bigl(g\sum_{j=1}^k(B_j-p_j)\bigr)$,
where $B_1,\ldots,B_k$ are independent with $B_j\sim\Bern(p_j)$ and
$0<p_j<1$. Then $(\mu_1,\ldots,\mu_n)\in\mathcal F_n(\mu)$ if and only if
there exists a partition of $\{1,\ldots,k\}$ into $n$ (possibly empty)
disjoint subsets $A_1,\ldots,A_n$ such that
\begin{equation}\label{eq:pb-factor-form}
 \mu_i=\Law\left(g\sum_{j\in A_i}(B_j-p_j)\right)\qquad(1\le i\le n),
\end{equation}
where an empty sum denotes $0$. If $n\ge k$, then $\mathcal M_n(\mu)$ is
equal to
\begin{equation}\label{eq:pb-min-price}
 \mathcal M_n(\mu)=g^3\sum_{j=1}^k p_j(1-p_j)\bigl(p_j^2+(1-p_j)^2\bigr),
\end{equation}
and the minimum is attained if and only if $(\mu_1,\ldots,\mu_n)$ is a
permutation of the $k$ singleton two-point laws $\Law(g(B_j-p_j))$ ($1\le j\le k$)
together with $n-k$ Dirac masses at $0$.

\item\label{it:binomial-equivalence} In particular, for any $n$-summand
independent array $(X_1,\ldots,X_n)$ that attains the $n$-summand supremum
$C_n$, the following two statements are equivalent:
\begin{enumerate}[label=(\alph*),leftmargin=2em]
\item $(X_1,\ldots,X_n)$ consists of $m\in\{1,\ldots,n\}$ identically
distributed nondegenerate two-point summands and $n-m$ zero summands;
\item the marginal distribution of the sum $S=X_1+\cdots+X_n$ is a centered
binomial distribution $\Law(g(\Bin(m,p)-mp))$ for some integer
$m\in\{1,\ldots,n\}$, some $g>0$, and some $p\in(0,1)$.
\end{enumerate}

\item\label{it:fair-binomial-oversized} For each integer $N\ge1$, let
$\mu_N=\Law\bigl((\Bin(N,1/2)-N/2)/\sqrt{N/4}\bigr)$ denote the standardized
symmetric binomial distribution of order $N$. Then for every capacity $n\ge1$,
every oversized symmetric binomial sum law ($N>n$) has strictly smaller
optimized ratio $\mathcal Q_n(\mu_N)=\Delta(\mu_N)/\mathcal M_n(\mu_N)$ than
$\mu_n$, and
\begin{equation}\label{eq:fair-binomial-sup}
 \sup_{N\ge1}\mathcal Q_n(\mu_N)
 =\begin{cases}
 R_1(1/2)=\Phi(1)-1/2 & \text{if }n=1,\\[3pt]
 \max\bigl\{R_n(1/2),\,R_{n-1}(1/2)\bigr\} & \text{if }n\ge2.
 \end{cases}
\end{equation}
\end{enumerate}
\end{theorem}

\begin{proof}
\ref{it:pb-fiber} Let $(X_1,\ldots,X_n)$ be independent centered random
variables with $\Law(X_i)=\mu_i$ and $\Law(X_1+\cdots+X_n)=\mu$. The support
of $\mu$ is the finite arithmetic progression
$-g\sum_{j=1}^k p_j+g\{0,1,\ldots,k\}$. Fixing one point in $\supp(X_r)$ for
each $r\ne i$, the sum of any point in $\supp(X_i)$ with those fixed points
lies in $\supp(\mu)$ because the $X_r$ are independent. Consequently,
$\supp(X_i)$ is finite and all differences between points of $\supp(X_i)$
belong to $g\Z$. Let $a_i=\min\supp(X_i)$; then $Y_i=(X_i-a_i)/g$ takes
nonnegative integer values with $\Pp(Y_i=0)>0$, and independence gives
$\sum_{i=1}^n a_i=\min\supp(\mu)=-g\sum_{j=1}^k p_j$ and
$\Law(Y_1+\cdots+Y_n)=\Law(B_1+\cdots+B_k)$.

Let $P_i(z)=\E z^{Y_i}$ ($1\le i\le n$) be the probability generating
polynomials of $Y_1,\ldots,Y_n$. Then
\[
 \prod_{i=1}^n P_i(z)=\prod_{j=1}^k(1-p_j+p_j z)\qquad(z\in\mathbb C).
\]
By unique factorization in $\mathbb C[z]$ and the normalization $P_i(1)=1$,
each polynomial $P_i(z)$ is the product of $(1-p_j+p_j z)$ over a subset
$A_i\subseteq\{1,\ldots,k\}$, where $A_1,\ldots,A_n$ form a partition of
$\{1,\ldots,k\}$. Since $\E X_i=0$ determines the translation $a_i=-g\sum_{j\in A_i}p_j$,
the representation \eqref{eq:pb-factor-form} follows.

Now suppose $n\ge k$. By repeated application of
Proposition~\ref{prop:splitting}, for every subset $A_i\subseteq\{1,\ldots,k\}$,
\[
 \E\left|g\sum_{j\in A_i}(B_j-p_j)\right|^3
 \ge g^3\sum_{j\in A_i}\E|B_j-p_j|^3,
\]
with equality if and only if $|A_i|\le1$. Summing over $i=1,\ldots,n$ yields
\eqref{eq:pb-min-price}, with equality if and only if every $A_i$ has
cardinality $0$ or $1$---that is, $k$ of the subsets $A_i$ are singletons and
the remaining $n-k$ are empty.

\ref{it:binomial-equivalence} Implication (a)$\implies$(b) is immediate from
the definition of the binomial distribution. Conversely, suppose
$(X_1,\ldots,X_n)$ attains $C_n$ and its sum $S$ has distribution
$\mu=\Law(g(\Bin(m,p)-mp))$ with $1\le m\le n$, $g>0$, and $0<p<1$. By
\eqref{eq:maximizer-minimizes-price}, the summand laws $\mu_i=\Law(X_i)$ must
achieve the minimum $\mathcal M_n(\mu)$ on $\mathcal F_n(\mu)$. Applying
part~\ref{it:pb-fiber} with $k=m$ and $p_1=\cdots=p_m=p$ shows that the
unique minimizers of $\sum_{i=1}^n\rho(\mu_i)$ on $\mathcal F_n(\mu)$ up to
coordinate permutation are $m$ copies of $\Law(g(B_1-p))$ and $n-m$ copies of
$\delta_0$, which is (a).

\ref{it:fair-binomial-oversized} By \citet[Theorem~1.3(a)]{MattnerSchulz2018},
for the standardized symmetric binomial law $\mu_N$ ($p=1/2$), the maximum
Kolmogorov discrepancy $d_N:=\Delta(\mu_N)$ is attained at the central
thresholds $\lfloor N/2\rfloor$ (on the right-continuous branch) and
$\lceil N/2\rceil$ (on the left-continuous branch), and equals
\begin{equation}\label{eq:dn-exact}
 d_{2k}=\frac12\binom{2k}{k}2^{-2k},\qquad
 d_{2k+1}=\Phi\!\left(\frac{1}{\sqrt{2k+1}}\right)-\frac12
 =\int_0^{1/\sqrt{2k+1}}\phi(u)\,du.
\end{equation}
We first verify that $d_N$ is strictly decreasing in $N\ge1$:
$d_1>d_2>d_3>\cdots$. Writing $u_k=\binom{2k}{k}2^{-2k}$ and
$I_r=\int_0^{\pi/2}\sin^r\theta\,d\theta$, integration by parts gives
$I_{2k}=(\pi/2)u_k$, $I_{2k+1}=1/\bigl((2k+1)u_k\bigr)$, and
$I_{2k-1}=1/(2ku_k)$; since $I_{2k+1}<I_{2k}<I_{2k-1}$ for all $k\ge1$,
Wallis's inequalities
\[
 \frac{1}{\sqrt{\pi(k+1/2)}}<u_k<\frac{1}{\sqrt{\pi k}}
\]
hold for all $k\ge1$. Since $\phz=1/\sqrt{2\pi}$, the left inequality gives
\[
 d_{2k}=\frac{u_k}{2}>\frac{\phz}{\sqrt{2k+1}}
 >\int_0^{1/\sqrt{2k+1}}\phi(u)\,du=d_{2k+1}.
\]
On the other hand, for every odd $r=2k-1\ge1$, the inequality
$e^{-u^2/2}\ge 1-u^2/2$ and the upper Wallis bound give
\[
 d_r=\int_0^{1/\sqrt r}\phi(u)\,du
 \ge\frac{\phz}{\sqrt r}\left(1-\frac{1}{6r}\right)
 >\frac{\phz}{\sqrt{r+1}}>\frac{u_k}{2}=d_{r+1},
\]
because $(1+1/r)(1-1/(6r))^2-1=(24r^2-11r+1)/(36r^3)>0$ for $r\ge1$. Thus
$d_N>d_{N+1}$ for all $N\ge1$.

Now let $N>n$. By \eqref{eq:mn-bounds}, $\mathcal M_n(\mu_N)\ge 1/\sqrt n$,
whereas for $N=n$ part~\ref{it:pb-fiber} gives $\mathcal M_n(\mu_n)=1/\sqrt n$.
Since $d_N<d_n$, we obtain
\[
 \mathcal Q_n(\mu_N)=\frac{d_N}{\mathcal M_n(\mu_N)}
 \le\sqrt n\,d_N<\sqrt n\,d_n=\frac{d_n}{\mathcal M_n(\mu_n)}=R_n(1/2).
\]
For $1\le N\le n$, part~\ref{it:pb-fiber} gives $\mathcal M_n(\mu_N)=1/\sqrt N$,
so $\mathcal Q_n(\mu_N)=\sqrt N\,d_N=R_N(1/2)$. When $n=1$ this leaves only
$R_1(1/2)$. When $n\ge2$, within each parity $R_N(1/2)$ is strictly
increasing in $N$: for odd $N$,
$R_N(1/2)=\sqrt N\int_0^{1/\sqrt N}\phi(u)\,du=\int_0^1\phi(v/\sqrt N)\,dv$
increases strictly with $N$, and for even $N=2k$,
$R_{2k+2}(1/2)/R_{2k}(1/2)=(2k+1)/\bigl(2\sqrt{k(k+1)}\bigr)>1$. Therefore
$\max_{1\le N\le n}R_N(1/2)=\max\{R_n(1/2),R_{n-1}(1/2)\}$, proving
\eqref{eq:fair-binomial-sup}.
\end{proof}

\begin{remark}[Balanced groupings for general $N>n$, and Gaussian and Poisson fibers]\label{rem:balanced-and-id}
When $\mu=\Law(g(\Bin(N,p)-Np))$ with $N>n$ and $0<p<1$,
Theorem~\ref{thm:binomial-fiber}\ref{it:pb-fiber} shows that every
factorization in $\mathcal F_n(\mu)$ has the form
$\mu_i=\Law(g(\Bin(k_i,p)-k_ip))$ with nonnegative integers
$k_1+\cdots+k_n=N$. Writing $c_p(k)=\E|\Bin(k,p)-kp|^3$ and $U\sim\Bern(p)-p$,
the identity $\frac{d^2}{dx^2}\bigl(\E|x+U|^3-|x|^3\bigr)=6\bigl(\E|x+U|-|x|\bigr)\ge0$
and conditional Jensen's inequality imply that $k\mapsto c_p(k+1)-c_p(k)$ is
nondecreasing, so $c_p$ is convex on $\Z_{\ge0}$ and $\mathcal M_n(\mu)$ is
attained by balanced group sizes $k_i\in\{\lfloor N/n\rfloor,\lceil N/n\rceil\}$.

By contrast, for infinitely divisible sum laws $\mu$, the fiber minimum
$\mathcal M_n(\mu)$ behaves quite differently from the finite-lattice case:
\begin{enumerate}[label=(\arabic*),leftmargin=2em]
\item If $\mu=N(0,V)$, then by Cram\'er's theorem \citep{Cramer1936} every
factor in $\mathcal F_n(\mu)$ is normal $N(0,v_i)$ with $\sum_{i=1}^n v_i=V$,
and by strict convexity of $v\mapsto v^{3/2}$ the unique minimizer of
$\sum_{i=1}^n\rho(\mu_i)=\rho(N(0,1))\sum_{i=1}^n v_i^{3/2}$ consists of $n$
identically distributed $N(0,V/n)$ summands (with $\mathcal M_n(N(0,V))=\rho(N(0,V))/\sqrt n$).
\item If $\mu=\Law(N_T-T)$ with $N_T\sim\Pois(T)$ ($T>0$), then for any
factorization $(X_1,\ldots,X_n)$ of $N_T-T$ the signed third moments add
($\sum_{i=1}^n\E X_i^3=\E(N_T-T)^3=T$), so
\[
 \sum_{i=1}^n\E|X_i|^3-T
 =\sum_{i=1}^n\bigl(\E|X_i|^3-\E X_i^3\bigr)
 =2\sum_{i=1}^n\E(X_i)_-^3>0,
\]
because every nonconstant centered variable takes strictly negative values
with positive probability. On the other hand, taking $n$ independent copies
of $N_{T/n}-T/n$ gives, for $n\ge T$,
\[
 \mathcal M_n(\mu)\le n\,\E|N_{T/n}-T/n|^3
 =T+\frac{2T^3}{n^2}e^{-T/n}\longrightarrow T\qquad(n\to\infty),
\]
so $\mathcal M_n(\mu)>T$ for every finite $n$ while
$\inf_{n\ge1}\mathcal M_n(\mu)=T$. (By Raikov's theorem \citep{Raikov1938},
every factor in $\mathcal F_n(\mu)$ is a centered Poisson law
$N_{t_i}-t_i$ with $\sum_{i=1}^n t_i=T$.) Notice that the limiting quantity
$T=\int y^3\,\nu(dy)$ is the third moment of the L\'evy jump measure
$\nu=T\delta_1$ arising from $n\to\infty$ rare summands, whereas at any fixed
finite $n$ every independent factorization of $N_T-T$ has third-moment sum
strictly greater than $T$.
\end{enumerate}
\end{remark}

Theorem~\ref{thm:binomial-fiber}\ref{it:binomial-equivalence} clarifies what
the structural classification question (Question~\ref{q:structure}) actually
requires: because exact-fiber third-moment minimization forces every
third-moment-minimizing factorization of $g(\Bin(m,p)-mp)$ ($m\le n$) in
$\mathcal F_n(\mu)$ to consist of $m$ identically distributed two-point
summands and $n-m$ zeros, proving that an $n$-summand maximizer of $C_n$ has
the padded i.i.d.\ two-point form \eqref{eq:padded-binomial-array} is
equivalent to proving that the marginal distribution of its sum $S$ is a
centered binomial distribution of order $m\le n$.

\section{Moment, transport, and characteristic-function criteria for identical two-point summands}\label{sec:criteria}

\subsection{The separated-support moment inequality}

Why do identically distributed two-point summands emerge when one compares
Kolmogorov discrepancy with the third absolute moment sum $R$? In his 1956
construction, \citet[Lemma~1, pp.~162--163]{Esseen1956} proved that for any
purely atomic random variable $X$ whose support points are separated by at
least $h>0$,
\begin{equation}\label{eq:esseen-moment-ineq}
 3h\Var(X)+\bigl|\E(X-\E X)^3\bigr|
 \le (3+\sqrt{10})\,\E|X-\E X|^3,
\end{equation}
with equality for nonconstant $X$ if and only if $X-\E X$ is a two-point
random variable of support span $h$ and upper probability $\pe$ or $1-\pe$.
By decomposing a centered distribution into its size-biased positive and
negative parts, we obtain an exact two-parameter family of sum-of-squares
identities containing \eqref{eq:esseen-moment-ineq} that classifies all
two-point equality laws of arbitrary skewness.

\begin{theorem}[Separated-support moment inequality and equality cases]\label{thm:separated-moment}
\begin{enumerate}[label=(\roman*),leftmargin=2em]
\item\label{it:one-law-moment} Let $X$ be a nonconstant centered real random
variable with finite third absolute moment whose support $\supp(X)$ satisfies
$|x-y|\ge h>0$ for all distinct $x,y\in\supp(X)$. Set
\[
 v=\E X^2,\qquad
 \kappa=\E X^3,\qquad
 M_+=\E(X_+)^3,\qquad
 M_-=\E(X_-)^3,\qquad
 \rho=\E|X|^3=M_++M_-,
\]
where $X_+=\max(X,0)$ and $X_-=\max(-X,0)$. Then
\begin{equation}\label{eq:geometric-gap-bound}
 hv\le\bigl(\sqrt{M_+}+\sqrt{M_-}\bigr)^2=\rho+2\sqrt{M_+M_-},
\end{equation}
and for every $\lambda>0$ and $\nu\in\R$, setting $d=\sqrt{\lambda^2+\nu^2}$,
\begin{equation}\label{eq:moment-sos}
 (\lambda+d)\rho-\lambda hv-\nu\kappa
 =\lambda\Bigl(\bigl(\sqrt{M_+}+\sqrt{M_-}\bigr)^2-hv\Bigr)
 +\Bigl(\sqrt{d-\nu}\sqrt{M_+}-\sqrt{d+\nu}\sqrt{M_-}\Bigr)^2\ge0.
\end{equation}
Equality holds in \eqref{eq:moment-sos} if and only if $X$ is supported on
two points $\{-a,b\}$ with $a,b>0$, $a+b=h$, and upper atom probability
\begin{equation}\label{eq:equality-prob}
 p=\Pp(X=b)=\frac{a}{h}=\frac12\left(1-\frac{\nu}{\lambda+d}\right).
\end{equation}

\item\label{it:array-moment} Let $(X_i)_{i\in I}$ be a finite or countable
collection of independent centered real random variables with
$0<V=\sum_{i\in I}\E X_i^2<\infty$ and $R=\sum_{i\in I}\E|X_i|^3<\infty$, and
suppose that for every nonconstant $X_i$ the points of $\supp(X_i)$ are
separated by at least the same $h>0$. Set $K_3=\sum_{i\in I}\E X_i^3$. Then
for every $\lambda>0$ and $\nu\in\R$ with $d=\sqrt{\lambda^2+\nu^2}$,
\begin{equation}\label{eq:array-moment-ineq}
 (\lambda+d)R-\lambda hV-\nu K_3\ge0,
\end{equation}
with equality if and only if $(X_i)_{i\in I}$ consists of finitely many
identically distributed two-point random variables of support span $h$ and
upper atom probability \eqref{eq:equality-prob}, together with deterministic
zero variables.
\end{enumerate}
\end{theorem}

\begin{proof}
\ref{it:one-law-moment} Because $X$ is centered and nonconstant,
$m:=\E X_+=\E X_->0$. Define positive random variables $B$ and $A$ with
distributions
\[
 \Pp(B\in E)=\frac{1}{m}\E\bigl[X_+\ind_{\{X\in E\cap(0,\infty)\}}\bigr],\qquad
 \Pp(A\in E)=\frac{1}{m}\E\bigl[X_-\ind_{\{-X\in E\cap(0,\infty)\}}\bigr].
\]
Then $v=\E X_+^2+\E X_-^2=m(\E B+\E A)$, $M_+=m\,\E B^2$, and
$M_-=m\,\E A^2$. For every $b\in\supp(B)\subset\supp(X)\cap(0,\infty)$ and
$a\in\supp(A)\subset(-\supp(X))\cap(0,\infty)$, the points $b>0$ and $-a<0$
are distinct elements of $\supp(X)$, so the separation hypothesis gives
$a+b\ge h$. Hence $\E A+\E B\ge h$, and the Cauchy--Schwarz inequality
$\E B\le\sqrt{\E B^2}$, $\E A\le\sqrt{\E A^2}$ yields
\[
 hv\le m(\E A+\E B)^2
 \le m\Bigl(\sqrt{\E A^2}+\sqrt{\E B^2}\Bigr)^2
 =\bigl(\sqrt{M_-}+\sqrt{M_+}\bigr)^2,
\]
which is \eqref{eq:geometric-gap-bound}. Expanding the two terms on the
right-hand side of \eqref{eq:moment-sos} and using $\rho=M_++M_-$,
$\kappa=M_+-M_-$, and $\sqrt{(d-\nu)(d+\nu)}=\lambda$ gives the identity
\eqref{eq:moment-sos}.

Both terms on the right-hand side of \eqref{eq:moment-sos} are nonnegative.
The first term vanishes if and only if $\E A+\E B=h$, $\Var(A)=0$, and
$\Var(B)=0$, which means that $A=a>0$ and $B=b>0$ are deterministic constants
with $a+b=h$. In that case $\supp(X)\setminus\{0\}=\{-a,b\}$; since $0<a,b<h$,
the point $0$ has distance $\min(a,b)<h$ to $\{-a,b\}$ and therefore cannot
belong to $\supp(X)$ under the hypothesis that all points of $\supp(X)$ are
separated by at least $h$. Thus $\supp(X)=\{-a,b\}$, and $\E X=0$ forces
$p=\Pp(X=b)=a/h$ and $q=\Pp(X=-a)=b/h$, with $M_+=p b^3=h^3pq^3$ and
$M_-=q a^3=h^3p^3q$. The second term on the right-hand side of
\eqref{eq:moment-sos} vanishes if and only if
\[
 \frac{q}{p}=\frac{b}{a}=\sqrt{\frac{M_+}{M_-}}
 =\sqrt{\frac{d+\nu}{d-\nu}}=\frac{d+\nu}{\lambda},
\]
which is equivalent to $q-p=(d+\nu-\lambda)/(d+\nu+\lambda)=\nu/(\lambda+d)$,
giving \eqref{eq:equality-prob}.

\ref{it:array-moment} Summing \eqref{eq:moment-sos} over all nonconstant
summands $X_i$ (while zero summands contribute $0$ to $R$, $V$, and $K_3$)
gives \eqref{eq:array-moment-ineq}. Equality holds if and only if every
nonconstant summand achieves equality in \eqref{eq:moment-sos} with the same
parameters $(h,\lambda,\nu)$, and hence has the same two-point distribution
of span $h$ and upper atom probability \eqref{eq:equality-prob}. Since each
such nonconstant summand has variance $h^2pq>0$ and $V=\sum_{i\in I}\E X_i^2<\infty$,
only finitely many summands can be nonconstant.
\end{proof}

\begin{remark}[Classical specializations]\label{rem:moment-specializations}
Setting $\lambda=3$ and $\nu=\pm1$ in Theorem~\ref{thm:separated-moment} gives
$d=\sqrt{10}$, $\lambda+d=3+\sqrt{10}$, and upper atom probability
$p=\frac12\bigl(1\mp 1/(3+\sqrt{10})\bigr)\in\{\pe,1-\pe\}$, recovering
Esseen's inequality \eqref{eq:esseen-moment-ineq} (which governs the leading
Edgeworth--Euler--Maclaurin term for half-line probabilities $\Pp(S\le x)$).
Setting $\lambda=1$ and $\nu=0$ gives $d=1$, $\lambda+d=2$, and $hv\le2\rho$,
with equality uniquely at symmetric two-point laws ($p=1/2$); this is the
moment inequality governing the asymptotic Berry--Esseen constant
$\sqrt{2/\pi}$ for interval probabilities $\Pp(a<S\le b)$ established by
\citet[Theorem~2.1 and Corollary~2.3]{DinevMattner2013}. Note also that the
assumption in Theorem~\ref{thm:separated-moment}\ref{it:one-law-moment} that
\emph{all} support points of $X$ are separated by at least $h$ (and not
merely that opposite-signed support points satisfy $a+b\ge h$) is necessary
for the two-point equality conclusion: if an atom at $0$ with probability
$r\in(0,1)$ is permitted at distance less than $h$ from $\{-a,b\}$, then the
centered three-point distribution assigning probabilities $(1-r)b/h$ to $-a$,
$r$ to $0$, and $(1-r)a/h$ to $b$ (where $a,b>0$, $a+b=h$, and
$b/a=(d+\nu)/\lambda$) also makes the right-hand side of
\eqref{eq:moment-sos} vanish.
\end{remark}

\subsection{An endogenous characteristic-function criterion}

Theorem~\ref{thm:separated-moment} assumes a common lower bound $h>0$ on the
support spacings of the summands. Remarkably, the moments of an arbitrary
independent array determine a canonical scale $\widehat h$ and a single
characteristic-function value whose modulus equals $1$ if and only if the
array is a finite i.i.d.\ two-point core with zero padding---without any a
priori discreteness, lattice, or support-separation assumption.

\begin{theorem}[Characteristic-function criterion and moment capacity bound]\label{thm:fourier-criterion}
Let $(X_i)_{i\in I}$ be a finite or countable collection of independent
centered real random variables with $0<V=\sum_{i\in I}\E X_i^2<\infty$ and
$R=\sum_{i\in I}\E|X_i|^3<\infty$, and let $S=\sum_{i\in I}X_i$ (converging
in $L^2$ and almost surely). Set
\begin{equation}\label{eq:moment-aggregates}
\begin{aligned}
 A&=\sum_{i\in I}\E(X_{i,+})^3,&
 B&=\sum_{i\in I}\E(X_{i,-})^3,\\
 K_3&=A-B=\sum_{i\in I}\E X_i^3,\qquad&
 D&=2\sqrt{AB}=\sqrt{R^2-K_3^2}>0,
\end{aligned}
\end{equation}
and define the moment-determined span, skewness, effective active count, and
characteristic-function deficit by
\begin{equation}\label{eq:endogenous-quantities}
\begin{aligned}
 \widehat h&=\frac{R+D}{V},&
 \widehat s&=\frac{K_3}{R+D}\in(-1,1),\\
 \widehat m&=\frac{2V^3}{D(R+D)},\qquad&
 \mathcal J&=1-\left|\E\exp\!\left(\frac{2\pi\ii S}{\widehat h}\right)\right|^2\ge0.
\end{aligned}
\end{equation}
Then:
\begin{enumerate}[label=(\roman*),leftmargin=2em]
\item\label{it:mhat-bound} If the index set $I=\{1,\ldots,n\}$ is finite, then
\begin{equation}\label{eq:mhat-le-n}
 0<\widehat m\le n,
\end{equation}
with equality $\widehat m=n$ if and only if all $n$ summands $X_1,\ldots,X_n$
are identically distributed nondegenerate two-point random variables.

\item\label{it:j-zero} For any finite or countable index set $I$,
\begin{equation}\label{eq:j-zero-equiv}
 \mathcal J=0
 \quad\Longleftrightarrow\quad
 \begin{aligned}[t]
 &(X_i)_{i\in I}\text{ consists of finitely many i.i.d.}\\
 &\text{nondegenerate two-point summands and zeros.}
 \end{aligned}
\end{equation}
When $\mathcal J=0$, the number of nondegenerate summands equals
$\widehat m\in\N$, their common support diameter is $\widehat h$, and their
common upper atom probability is $p=(1-\widehat s)/2$.

\item\label{it:j-stability} Fix a finite horizon $n\ge1$ and $L_0<\infty$. If
$(X_1^{(k)},\ldots,X_n^{(k)})$ is a sequence of $n$-summand independent
centered arrays with $V^{(k)}=1$, $R^{(k)}\le L_0$, and
$\mathcal J^{(k)}\to0$ as $k\to\infty$, then every subsequence has a further
subsequence converging coordinatewise in $W_3$ to an array consisting of
$m\in\{1,\ldots,n\}$ identically distributed nondegenerate two-point summands
and $n-m$ zeros.
\end{enumerate}
\end{theorem}

\begin{proof}
First note that since $V>0$ and $\E X_i=0$, at least one $X_i$ is nonconstant
and takes both positive and negative values with positive probability, so
$A>0$, $B>0$, and $D=2\sqrt{AB}>0$.

\ref{it:mhat-bound} For a single nonconstant centered summand $X_i$, write
$v_i=\E X_i^2$, $A_i=\E(X_{i,+})^3$, $B_i=\E(X_{i,-})^3$,
$\rho_i=A_i+B_i$, $\delta_i=2\sqrt{A_iB_i}$, and $m_i=\E X_{i,+}=\E X_{i,-}>0$.
By H\"older's inequality on $\{X_i>0\}$ and $\{X_i<0\}$,
\[
 m_i=\E\bigl[X_{i,+}\ind_{\{X_i>0\}}\bigr]\le A_i^{1/3}\Pp(X_i>0)^{2/3},\qquad
 m_i=\E\bigl[X_{i,-}\ind_{\{X_i<0\}}\bigr]\le B_i^{1/3}\Pp(X_i<0)^{2/3}.
\]
Since $\Pp(X_i>0)+\Pp(X_i<0)\le1$, summing the $3/2$-powers gives
\[
 m_i^{3/2}\left(\frac{1}{\sqrt{A_i}}+\frac{1}{\sqrt{B_i}}\right)\le 1,
 \qquad\text{or}\qquad
 m_i\le\left(\frac{\sqrt{A_iB_i}}{\sqrt{A_i}+\sqrt{B_i}}\right)^{2/3}.
\]
Meanwhile, Cauchy--Schwarz gives
$v_i=\E X_{i,+}^2+\E X_{i,-}^2\le\sqrt{m_i}\bigl(\sqrt{A_i}+\sqrt{B_i}\bigr)$.
Raising to the power $3/2$ and substituting the bound on $m_i$ yields
\begin{equation}\label{eq:vi-three-halves}
 v_i^{3/2}\le m_i^{3/4}\bigl(\sqrt{A_i}+\sqrt{B_i}\bigr)^{3/2}
 \le\sqrt{\sqrt{A_iB_i}}\,\bigl(\sqrt{A_i}+\sqrt{B_i}\bigr)
 =\sqrt{\frac{\delta_i(\rho_i+\delta_i)}{2}},
\end{equation}
which also holds trivially when $X_i=0$. Equality holds in
\eqref{eq:vi-three-halves} for $v_i>0$ if and only if $\Pp(X_i=0)=0$ and
$X_{i,+}$, $X_{i,-}$ are constant on $\{X_i>0\}$ and $\{X_i<0\}$, meaning
that $X_i$ is a two-point random variable. By Cauchy--Schwarz and Minkowski's
inequality $\sum_{i=1}^n\delta_i=2\sum_{i=1}^n\sqrt{A_iB_i}\le 2\sqrt{AB}=D$,
summing \eqref{eq:vi-three-halves} over $i=1,\ldots,n$ gives
\[
 \frac{V^{3/2}}{\sqrt n}
 \le\sum_{i=1}^n v_i^{3/2}
 \le\sqrt{\frac12\left(\sum_{i=1}^n\delta_i\right)\left(R+\sum_{i=1}^n\delta_i\right)}
 \le\sqrt{\frac{D(R+D)}{2}},
\]
which rearranges to $\widehat m=2V^3/\bigl(D(R+D)\bigr)\le n$. Equality
$\widehat m=n$ requires $v_1=\cdots=v_n=V/n>0$, each $X_i$ two-point, and
$(A_i,B_i)$ proportional across all $i$, which forces all $n$ two-point
summands to be identically distributed.

\ref{it:j-zero} If $(X_i)_{i\in I}$ consists of $m\ge1$ identically
distributed two-point summands $g(B_i-p)$ ($g>0$, $0<p<1$, $q=1-p$) and zeros,
then $V=mg^2pq$, $A=mg^3pq^3$, $B=mg^3p^3q$, $R=mg^3pq(p^2+q^2)$,
$K_3=mg^3pq(q-p)$, and $D=2mg^3p^2q^2$, so $R+D=mg^3pq(p+q)^2=mg^3pq$. Hence
$\widehat h=(R+D)/V=g$, $\widehat s=q-p$, $\widehat m=m$, and since every
summand is supported on a coset of $g\Z$,
$\bigl|\E e^{2\pi\ii S/g}\bigr|=\prod_{i\in I}\bigl|\E e^{2\pi\ii X_i/g}\bigr|=1$,
so $\mathcal J=0$.

Conversely, suppose $\mathcal J=0$. By independence,
\[
 1=\left|\E\exp\!\left(\frac{2\pi\ii S}{\widehat h}\right)\right|^2
 =\prod_{i\in I}\left|\E\exp\!\left(\frac{2\pi\ii X_i}{\widehat h}\right)\right|^2,
\]
so $\bigl|\E\exp(2\pi\ii X_i/\widehat h)\bigr|=1$ for every $i\in I$. Thus
for each $i\in I$, $\exp(2\pi\ii X_i/\widehat h)$ is almost surely constant,
so $\supp(X_i)$ is contained in a single coset of $\widehat h\Z$. In
particular, for every nonconstant $X_i$, the points of $\supp(X_i)$ are
separated by at least $\widehat h$. Applying
Theorem~\ref{thm:separated-moment}\ref{it:array-moment} with span
$h=\widehat h$, $\lambda=1$, and $\nu=K_3/D$, we have
$d=\sqrt{1+K_3^2/D^2}=R/D$ and $\lambda+d=(R+D)/D$, so the left-hand side of
\eqref{eq:array-moment-ineq} equals
\[
 \frac{R+D}{D}R-\frac{R+D}{V}V-\frac{K_3}{D}K_3
 =\frac{R(R+D)-K_3^2}{D}-(R+D)
 =\frac{R^2-K_3^2+RD}{D}-(R+D)=0,
\]
because $R^2-K_3^2=D^2$. Therefore equality holds in
\eqref{eq:array-moment-ineq}, and Theorem~\ref{thm:separated-moment}\ref{it:array-moment}
implies that $(X_i)_{i\in I}$ consists of finitely many identically
distributed two-point summands of span $\widehat h$ and upper probability
$p=\frac12(1-\nu/(\lambda+d))=(1-\widehat s)/2$, together with zeros.

\ref{it:j-stability} Let $(X_1^{(k)},\ldots,X_n^{(k)})$ have $V^{(k)}=1$,
$R^{(k)}\le L_0$, and $\mathcal J^{(k)}\to0$. Because
$\sum_{i=1}^n\E|X_i^{(k)}|^3\le L_0$, the sequence of $n$-tuples is tight in
$\R^n$ with uniformly integrable second moments. Passing to a subsequence, we
may assume that $X_i^{(k)}\Rightarrow X_i^{(\infty)}$ weakly with
$\E X_i^{(\infty)}=0$ and $\sum_{i=1}^n\E(X_i^{(\infty)})^2=1$, while
$A_i^{(k)}=\E(X_{i,+}^{(k)})^3\to a_i$ and
$B_i^{(k)}=\E(X_{i,-}^{(k)})^3\to b_i$. By Fatou's lemma for weak convergence,
$A_i':=\E(X_{i,+}^{(\infty)})^3\le a_i$ and
$B_i':=\E(X_{i,-}^{(\infty)})^3\le b_i$. Let $A'=\sum_{i=1}^n A_i'>0$,
$B'=\sum_{i=1}^n B_i'>0$, $a=\sum_{i=1}^n a_i\ge A'$, and
$b=\sum_{i=1}^n b_i\ge B'$. Then
\[
 \widehat h^{(k)}=\bigl(\sqrt{\textstyle\sum_i A_i^{(k)}}+\sqrt{\textstyle\sum_i B_i^{(k)}}\bigr)^2
 \longrightarrow H:=\bigl(\sqrt a+\sqrt b\bigr)^2\in\bigl[n^{-1/2},2L_0\bigr].
\]
Since $|\E e^{\ii t X_i^{(k)}}-\E e^{\ii u X_i^{(k)}}|\le|t-u|\,\E|X_i^{(k)}|\le|t-u|$,
weak convergence $X_i^{(k)}\Rightarrow X_i^{(\infty)}$ and
$\widehat h^{(k)}\to H>0$ imply
$|\E e^{2\pi\ii S^{(\infty)}/H}|^2=\lim_{k\to\infty}(1-\mathcal J^{(k)})=1$.
Thus each nonconstant limiting summand $X_i^{(\infty)}$ is supported on a
coset of $H\Z$ and has support spacing at least $H$. Summing
\eqref{eq:geometric-gap-bound} of Theorem~\ref{thm:separated-moment} over the
nonconstant limiting summands (with total variance $1$) and using Minkowski's
inequality gives
\[
 H=H\sum_{i=1}^n\E(X_i^{(\infty)})^2
 \le\sum_{i=1}^n\Bigl(\sqrt{A_i'}+\sqrt{B_i'}\Bigr)^2
 \le\Bigl(\sqrt{A'}+\sqrt{B'}\Bigr)^2
 \le\Bigl(\sqrt a+\sqrt b\Bigr)^2=H.
\]
Hence every inequality in this chain is an equality: $A'=a$, $B'=b$, and so
$A_i'=a_i$ and $B_i'=b_i$ for each $i=1,\ldots,n$. Thus
$\E|X_i^{(k)}|^3\to\E|X_i^{(\infty)}|^3$ alongside weak convergence, which
implies $W_3(X_i^{(k)},X_i^{(\infty)})\to0$, and $\mathcal J^{(\infty)}=0$
together with part~\ref{it:j-zero} shows that $(X_1^{(\infty)},\ldots,X_n^{(\infty)})$
is a padded i.i.d.\ two-point array.
\end{proof}

\begin{remark}[Diagnostic equality criteria versus objective monotonicity]\label{rem:j-not-monotone}
Theorem~\ref{thm:fourier-criterion} characterizes the exact zero set
$\mathcal J=0$ across the entire domain of independent arrays, without any
prior restriction to discrete or common-diameter distributions. However,
$\mathcal J=0$ is a structural equality criterion rather than a monotonicity
statement for the Berry--Esseen ratio $\mathcal C$: neither
Theorem~\ref{thm:separated-moment} nor Theorem~\ref{thm:fourier-criterion}
proves that a global maximizer of $\mathcal C$ must satisfy $\mathcal J=0$,
or that decreasing $\mathcal J$ increases $\mathcal C$. As we discuss in
Section~\ref{sec:open}, proving the structural conjecture by this route
requires supplying the missing global implication---namely, showing that every
$n$-summand array with $\mathcal J>0$ admits an admissible replacement with
strictly larger Berry--Esseen ratio $\mathcal C$.
\end{remark}

\subsection{Comparison with zero-bias Wasserstein transport}\label{subsec:transport-comparison}

A third global characterization of padded i.i.d.\ two-point arrays comes from
Stein's method and the zero-bias transformation \citep{GoldsteinReinert1997}.
For a centered random variable $X$ with variance $v=\E X^2>0$, its
\emph{zero-biased distribution} $X^*$ is characterized by the identity
$\E[Xf(X)]=v\,\E f'(X^*)$ for all absolutely continuous functions $f:\R\to\R$
with bounded derivative. For a finite independent array $(X_1,\ldots,X_n)$
with sum $S$, total variance $V>0$, and third-moment sum $R<\infty$, the
$L^1$-Wasserstein distance $W_1(S,S^*)=\int_\R|\Pp(S\le x)-\Pp(S^*\le x)|\,dx$
satisfies the normalized transport inequality
\begin{equation}\label{eq:transport-ineq}
 \mathcal B(X_1,\ldots,X_n):=\frac{2VW_1(S,S^*)}{R}\le 1
\end{equation}
\citep{Goldstein2010,Tyurin2009}. In the companion paper
\citet[Theorem~1.1]{BellTransport2026}, it is proved via Choquet's theorem on
fixed-variance moment sets and a conditional-variance identity that
\begin{equation}\label{eq:transport-rigidity}
 \mathcal B(X_1,\ldots,X_n)=1
 \quad\Longleftrightarrow\quad
 \begin{aligned}[t]
 &(X_1,\ldots,X_n)\text{ consists of i.i.d.}\\
 &\text{nondegenerate two-point summands and zeros.}
 \end{aligned}
\end{equation}
Furthermore, let $S_1,S_2,\ldots$ be independent copies of the block sum
$S=X_1+\cdots+X_n$, and define the $m$-fold repeated-block Berry--Esseen
ratio (which is distinct from the unrestricted $m$-summand supremum $C_m$) by
\begin{equation}\label{eq:repeated-block-def}
\begin{aligned}
 \mathcal C_{\mathrm{rep},m}(X_1,\ldots,X_n)
 &:=\frac{(mV)^{3/2}\Delta(S_1+\cdots+S_m)}{mR}\\
 &=\frac{\sqrt m\,V^{3/2}}{R}\sup_{x\in\R}\left|\Pp\left\{\frac{S_1+\cdots+S_m}{\sqrt{mV}}\le x\right\}-\Phi(x)\right|.
\end{aligned}
\end{equation}
Combining \eqref{eq:transport-rigidity} with the sharp lattice transport
inequality
\[
 |\E X^3|+3h\Var(X)\le(3+\sqrt{10})\,2\Var(X)W_1(X,X^*)
\]
for variables supported on a coset of $h\Z$ proves that for any fixed finite
block $(X_1,\ldots,X_n)$,
\begin{equation}\label{eq:repeated-block-bound}
 \lim_{m\to\infty}\mathcal C_{\mathrm{rep},m}(X_1,\ldots,X_n)
 \le\ce\,\mathcal B(X_1,\ldots,X_n)\le\ce,
\end{equation}
with equality in \eqref{eq:repeated-block-bound} if and only if all
nonconstant summands in the block share a single Esseen two-point law with
probability $\pe$ or $1-\pe$ \citep[Theorem~1.2]{BellTransport2026}. Along any
sequence of independent arrays with $L_n\to0$ and $\mathcal C_n\to\ce$, one
likewise has $\mathcal B_n\to1$ \citep[Theorem~1.1]{BellSmallStructure2026}.

One might ask whether the asymptotic bound \eqref{eq:repeated-block-bound}
holds pointwise for finite arrays---that is, whether
$\mathcal C(X_1,\ldots,X_n)\le\ce\,\mathcal B(X_1,\ldots,X_n)$ holds for all
independent arrays. The following explicit one-summand family shows that no
constant $K<\infty$ can satisfy $\mathcal C(X)\le K\,\mathcal B(X)$ for all
centered distributions $X$.

\begin{proposition}[Failure of pointwise domination by zero-bias transport]\label{prop:transport-counterexample}
For each $a>0$, let $X_a$ be the centered real random variable that agrees
with $Z\sim N(0,1)$ on $\{|x|>a\}$ and replaces the restriction of $Z$ to
$[-a,a]$ by three atoms at $\{-a,0,a\}$ having the same total mass, zero
mean, and second moment as $Z\ind_{\{|Z|\le a\}}$. Then $\E X_a=0$,
$\Var(X_a)=1$, and as $a\downarrow0$,
\begin{equation}\label{eq:xa-asymptotics}
 \mathcal C(X_a)=\frac{a}{6}+O(a^3),\qquad
 \mathcal B(X_a)=\frac{5a^2}{18}+O(a^4),\qquad
 \frac{\mathcal C(X_a)}{\mathcal B(X_a)}=\frac{3}{5a}+O(a)\longrightarrow\infty.
\end{equation}
\end{proposition}

\begin{proof}
Let $m_a=\int_{-a}^a\phi(u)\,du=2a\phz+O(a^3)$ and
$v_a=\int_{-a}^a u^2\phi(u)\,du=\frac23 a^3\phz+O(a^5)$. To match the mass,
zero mean, and second moment of $Z\ind_{\{|Z|\le a\}}$ on $\{-a,0,a\}$, $X_a$
places mass $w_a=v_a/(2a^2)=\frac13 a\phz+O(a^3)$ at $-a$ and at $+a$, and
mass $m_a-2w_a=\frac43 a\phz+O(a^3)$ at $0$. Then $\Var(X_a)=\Var(Z)=1$, and
\[
 R_a=\E|X_a|^3=\E|Z|^3+2a^3 w_a-\int_{-a}^a|u|^3\phi(u)\,du
 =4\phz+O(a^4).
\]
Outside $(-a,a)$, the distribution function $F_{X_a}(x)$ coincides with
$\Phi(x)$. On $(-a,0)$ it is constant and equal to $\Phi(-a)+w_a$, while on
$[0,a)$ it equals $\Phi(a)-w_a$. Thus the maximum of $|F_{X_a}(x)-\Phi(x)|$ is
attained as $x\uparrow0$ (and by symmetry at $x=0$):
\[
 \Delta(X_a)=\Phi(0)-\bigl(\Phi(-a)+w_a\bigr)
 =\frac12 m_a-w_a=\frac23 a\phz+O(a^3),
\]
giving $\mathcal C(X_a)=\Delta(X_a)/R_a=a/6+O(a^3)$.

On the other hand, because $\Var(X_a)=1$, the zero-biased distribution
$X_a^*$ has Lebesgue density $f_{X_a^*}(x)=\E[X_a\ind_{\{X_a>x\}}]$. For
$|x|>a$, $\E[X_a\ind_{\{|X_a|\le a\}}]=0$ implies
$f_{X_a^*}(x)=\int_x^\infty u\phi(u)\,du=\phi(x)$. For $x\in(-a,0)\cup(0,a)$,
\[
 f_{X_a^*}(x)=\int_a^\infty u\phi(u)\,du+a w_a=\phi(a)+\frac{v_a}{2a}.
\]
Consequently, $F_{X_a^*}(x)=\Phi(x)$ for $|x|\ge a$, and on $[-a,a]$ the
distribution function $F_{X_a^*}$ is affine with slope
$\phi(a)+v_a/(2a)=\phz+O(a^2)$, interpolating from $\Phi(-a)$ at $-a$ to
$\Phi(a)$ at $a$. Hence
\begin{align*}
 W_1(X_a,X_a^*)&=\int_{-a}^a\bigl|F_{X_a}(x)-F_{X_a^*}(x)\bigr|\,dx\\
 &=2\int_{-a}^0\bigl|w_a-(x+a)\bigl(\phz+O(a^2)\bigr)\bigr|\,dx\\
 &=2\phz\int_0^a\left|\frac{a}{3}-u\right|\,du+O(a^4)
 =2\phz\left(\frac{a^2}{18}+\frac{4a^2}{18}\right)+O(a^4)
 =\frac59 a^2\phz+O(a^4).
\end{align*}
Dividing $2W_1(X_a,X_a^*)$ by $R_a=4\phz+O(a^4)$ gives
$\mathcal B(X_a)=5a^2/18+O(a^4)$ and completes the proof of
\eqref{eq:xa-asymptotics}.
\end{proof}

\section{Known value bounds and partial classifications of maximizers}\label{sec:known}

\subsection{Attainment of $C_n$ and equivalence with positive-variance arrays}\label{subsec:attainment}

Before reviewing the known upper bounds and structural classifications, we
record the basic existence theorem for fixed-$n$ maximizers (see also
\citealp[Lemma~2.1]{BellSmallStructure2026}) and verify that permitting
deterministic zero summands does not enlarge the $n$-summand supremum $C_n$.

\begin{proposition}[Attainment of $C_n$ and positive-variance approximation]\label{prop:cn-attainment}
For every integer $n\ge1$:
\begin{enumerate}[label=(\roman*),leftmargin=2em]
\item\label{it:cn-attained} The supremum $C_n=\sup\mathcal C(X_1,\ldots,X_n)$
over $n$-tuples of independent centered real random variables with $V>0$ and
$R<\infty$ (permitting zero summands) is attained, and every maximizing
sequence normalized to $V=1$ has a subsequence converging coordinatewise in
$W_3$ to a maximizing array.
\item\label{it:cn-pos-var} The supremum of $\mathcal C(X_1,\ldots,X_n)$ over
$n$-tuples of independent centered real random variables satisfying
$\Var(X_i)>0$ and $\E|X_i|^3<\infty$ for every $i\in\{1,\ldots,n\}$ equals
$C_n$.
\end{enumerate}
\end{proposition}

\begin{proof}
\ref{it:cn-attained} Let $(X_1^{(k)},\ldots,X_n^{(k)})$ be a maximizing
sequence for $C_n$. By scale invariance we may assume
$V^{(k)}=\sum_{i=1}^n\E(X_i^{(k)})^2=1$. Since $C_n\ge C_1>0$ and
$\Delta(S^{(k)})\le1$, for all large $k$ we have
\[
 n^{-1/2}\le R^{(k)}=\frac{\Delta(S^{(k)})}{\mathcal C(X_1^{(k)},\ldots,X_n^{(k)})}
 \le\frac{2}{C_1}<\infty.
\]
For each coordinate $i\in\{1,\ldots,n\}$ and any $M>0$, Markov's inequality
gives
\[
 \E\bigl[(X_i^{(k)})^2\ind_{\{|X_i^{(k)}|>M\}}\bigr]
 \le\frac{R^{(k)}}{M}\le\frac{2}{MC_1}.
\]
Hence the laws of $X_i^{(k)}$ are tight with uniformly integrable first and
second moments. Passing to a subsequence, we may assume that
$X_i^{(k)}\Rightarrow X_i^{(\infty)}$ weakly as $k\to\infty$ for each
$i=1,\ldots,n$, with $\E X_i^{(\infty)}=0$,
$V^{(\infty)}=\sum_{i=1}^n\E(X_i^{(\infty)})^2=1$, and, by Fatou's lemma,
\[
 \E|X_i^{(\infty)}|^3\le\liminf_{k\to\infty}\E|X_i^{(k)}|^3\qquad(1\le i\le n).
\]
By independence, the sums $S^{(k)}=\sum_{i=1}^n X_i^{(k)}$ converge weakly and
in $W_1$ to $S^{(\infty)}=\sum_{i=1}^n X_i^{(\infty)}$. Because the normal
distribution function $\Phi$ is continuous with $\|\Phi'\|_\infty=\phz$, the
estimate $\bigl|\Delta(S^{(k)})-\Delta(S^{(\infty)})\bigr|\le\sqrt{2\phz W_1(S^{(k)},S^{(\infty)})}$
established in the proof of Theorem~\ref{thm:fiber-regularity}\ref{it:ratio-continuity}
(which includes both closed and strict traces at any atoms of $S^{(\infty)}$)
implies that $\Delta(S^{(k)})\to\Delta(S^{(\infty)})=C_n\lim_{k\to\infty}R^{(k)}>0$.
Consequently,
\[
 C_n=\frac{\Delta(S^{(\infty)})}{\lim_{k\to\infty}R^{(k)}}
 \le\frac{\Delta(S^{(\infty)})}{\sum_{i=1}^n\E|X_i^{(\infty)}|^3}\le C_n.
\]
Both inequalities must be equalities, so $(X_1^{(\infty)},\ldots,X_n^{(\infty)})$
attains $C_n$ and $\lim_{k\to\infty}\E|X_i^{(k)}|^3=\E|X_i^{(\infty)}|^3$ for
every $i=1,\ldots,n$, which together with weak convergence gives
$W_3(X_i^{(k)},X_i^{(\infty)})\to0$.

\ref{it:cn-pos-var} Let $(X_1,\ldots,X_n)$ be an array attaining $C_n$ with
$V=1$. For each $\eps>0$, define $X_i^{(\eps)}=X_i$ if $\Var(X_i)>0$ and let
$X_i^{(\eps)}\sim N(0,\eps)$ (independent of all other coordinates) if
$X_i=0$. Then $\Var(X_i^{(\eps)})>0$ for all $i=1,\ldots,n$, while as
$\eps\downarrow0$ we have $V^{(\eps)}\to 1$,
$R^{(\eps)}=R+O(\eps^{3/2})\to R$, and
$\Delta(S^{(\eps)})\to\Delta(S)$, so
$\mathcal C(X_1^{(\eps)},\ldots,X_n^{(\eps)})\to C_n$.
\end{proof}

Allowing zero summands is therefore the natural closed compactification of the
fixed-$n$ domain: as \citet[Theorem~1.2]{BellBernoulli2026} shows for $n=2$
two-point summands, the supremum over strictly positive-variance two-point
arrays equals $C_1$ because one summand's variance can degenerate to $0$, and
that supremum is attained only when the zero summand is included in the
domain.

\subsection{Quantitative upper bounds on $C_n$ and $\Cind$}\label{subsec:value-bounds}

We next summarize the quantitative upper bounds on $C_n$ and $\Cind$
established in the literature and in Manuscripts~1--3
\citep{BellIndependent2026,BellFinite2026,BellSmall2026}:

\begin{enumerate}[label=(\arabic*),leftmargin=2em]
\item \emph{One summand ($n=1$):} \citet[Theorem~1]{BentkusKirsa1989} proved
that for a single centered random variable $X_1$,
\begin{equation}\label{eq:c1-bound}
 C_1=B_1=\max_{0<p<1}\frac{\sqrt{p(1-p)}}{p^2+(1-p)^2}
 \left\{\Phi\!\left(\sqrt{\frac{p}{1-p}}\right)-p\right\}
 =0.370351523474378920\ldots<\ce,
\end{equation}
attained up to scale and sign reflection only by the two-point distribution at
$p_1=0.390958631281332815\ldots$.

\item \emph{Two arbitrary summands ($n=2$):} \citet[Theorem~1.1]{BellFinite2026}
proves that over all pairs of independent centered random variables,
\begin{equation}\label{eq:c2-bound}
 C_2<\frac{49}{120}=0.408333\ldots<\ce.
\end{equation}
The proof fixes the threshold and individual variances, uses strong quadratic
moment duality \citep{Winkler1988,BertsimasPopescu2005} to reduce each summand
to at most three support points, and rules out all combinatorial contact
configurations above $49/120$ by variational identities and certified interval
enclosures.

\item \emph{Two-point arrays and variance-concentrated arrays:} For $n$
identically distributed two-point summands, \citet[Theorem~1]{Schulz2016}
proved the strict binomial bound $B_n<\ce$. \citet[Theorems~1.2 and~1.3]{BellFinite2026}
extend the sharp inequality $\Delta(S)\le\ce L$ to two broader classes of
non-identically distributed two-point arrays:
\begin{enumerate}[label=(\alph*),leftmargin=2em]
\item arbitrary numbers $n\ge1$ of independent two-point summands whose
nondegenerate members share a common support diameter $h>0$;
\item at most three independent two-point summands ($n\le3$) with arbitrary
support diameters (with the uniform strict gap $\ce L-\Delta(S)>3/80$ when
$n\le2$).
\end{enumerate}
Combining these bounds with a Stein--Ornstein--Uhlenbeck normal comparison for
the remaining summands yields $\Delta(S)\le\ce L$ for any finite independent
array in which two two-point summands carry at least $22/23$ of the total
variance $V$, or two arbitrary summands carry at least $439/440$ of $V$
\citep[Corollary~1.4]{BellFinite2026}.

\item \emph{General independent arrays with small Lyapunov ratio:}
Adapting the large-sample i.i.d.\ method of \citet{HeCheng2026} to
non-identically distributed summands, \citet[Theorem~1.1]{BellSmall2026}
proves that for every finite independent array,
\begin{equation}\label{eq:small-lyapunov-value}
 L\le\exp\bigl(-\exp(97)\bigr)
 \quad\Longrightarrow\quad
 \Delta(S)\le\ce L.
\end{equation}
If a counterexample existed below this threshold, maximizing $\Delta/L$ minus
a smooth localization penalty (bounded using the asymptotic smoothing
inequality $\Delta\le\ce L+0.2538L^{4/3}$ of \citealp[Corollary~4.18]{Shevtsova2012})
would produce an attained extremal array with $|X_i|\le 8L\sqrt V$; uniform
jitter smoothing, quantitative stability in Esseen's moment inequality
\eqref{eq:esseen-moment-ineq}, and conditional Fourier and clipping
comparisons against a common-diameter two-point array (bounded by $\ce L$ via
\citealp[Theorem~1.2(ii)]{BellFinite2026}) then rule out any violation.

\item \emph{Unconditional upper bound across all independent arrays:}
\citet[Theorem~1.1]{BellIndependent2026} proves that for every $n\ge1$ and
every independent array,
\begin{equation}\label{eq:unconditional-upper}
 \Cind\le U_*<0.44988,
\end{equation}
improving upon the bound $0.5583$ of \citet{Shevtsova2013} by combining a
zero-bias characteristic-function differential inequality, a signed cosine-product
reference \citep{MattnerShevtsova2019}, isolation of the largest-variance
summand, and outward-rounded interval subdivision over the moment parameter
space.
\end{enumerate}

\subsection{General variational support reductions and compactification}\label{subsec:variational}

What can be proved about maximizing arrays without restricting either the
support of the summands or the size of the Lyapunov ratio $L$? In
Manuscript~7 \citep{BellExtremizers2026}, variational perturbations of a
single summand's distribution---fixing the maximizing threshold $z$ and the
deleted sum $T_i=\sum_{j\ne i}X_j$---establish two general necessary
conditions:

\begin{enumerate}[label=(\arabic*),leftmargin=2em]
\item \emph{Four-point and three-point support reduction}
\citep[Theorem~3.2]{BellExtremizers2026}: At any local maximum of $\mathcal C$
with respect to two-sided probability perturbations of $X_i$ preserving total
mass $1$ and mean $0$ (while allowing $v_i=\E X_i^2$ and $\rho_i=\E|X_i|^3$
to vary), the second logarithmic derivative of $\mathcal C$ restricts the
linear span of the moment vectors $(1,u,u^2,|u|^3)$ on $\supp(X_i)$, proving
that every summand of a maximizer is supported on at most four points. At a
global maximum, any four-point summand lies on a flat probability segment
whose endpoint replaces $X_i$ by a three-point summand with the exact same
variance $v_i$ and third absolute moment $\rho_i$ (and the same ratio
$\mathcal C$).

\item \emph{Convolution indecomposability when a spare slot is available}
\citep[Corollary~3.4]{BellExtremizers2026}: By the third-moment splitting
identity (Proposition~\ref{prop:splitting}), every nonconstant summand $X_i$
of an unrestricted-horizon maximizer---or of a fixed-$n$ maximizer that has at
least one zero summand---is convolution-indecomposable, since splitting a
decomposable summand into two nonconstant independent centered factors
preserves the sum $S$ and total variance $V$ while strictly decreasing $R$.
In particular, if such a maximizer has a three-point summand supported on an
arithmetic progression $\{x_0,x_0+g,x_0+2g\}$ ($g>0$) with probabilities
$(m_0,m_1,m_2)$, then $m_1^2<4m_0m_2$ (equivalently, $\Var(X_i)>ab/2$ when
the centered endpoints are $-a<0<b$), since otherwise its probability
generating polynomial factors into two nondegenerate two-point laws of
diameter $g$.
\end{enumerate}

When studying the universal constant $\Cind=\sup_{n\ge1}C_n$, if $\Cind>\ce$
then the small-Lyapunov bound \eqref{eq:small-lyapunov-value} forces every
unit-variance maximizing sequence to satisfy $R^{(k)}>\exp(-\exp(97))$.
Passing to the limit along variances ordered nonincreasingly produces an
attained maximizer in a L\'evy--Khintchine compactification
\citep[Definition~1.1 and Theorem~1.1]{BellExtremizers2026}:
\begin{equation}\label{eq:compactified-rep}
\begin{aligned}
 S&=\sum_{i\ge1}X_i+Y,\qquad
 \sum_{i\ge1}\Var(X_i)+a+\int_{\R\setminus\{0\}}y^2\,\nu(dy)=1,\\
 \beta&=\sum_{i\ge1}\E|X_i|^3+\int_{\R\setminus\{0\}}|y|^3\,\nu(dy),
\end{aligned}
\end{equation}
where $(X_i)_{i\ge1}$ are independent centered discrete summands (each
reducible to at most three support points), $Y$ is an independent centered
infinitely divisible random variable with Gaussian variance $a\ge0$ and
L\'evy measure $\nu$, and $\Delta(S)=\Cind\beta$. Notice that in
\eqref{eq:compactified-rep}, the Gaussian component $a$ contributes $0$ to
$\beta$ (as a limit of $N\to\infty$ summands of variance $a/N$ whose
third-moment sum $O(N^{-1/2})$ vanishes), while the jump measure $\nu$
contributes the limiting rare-summand sum $\int|y|^3\,\nu(dy)$. Even in this
enlarged compactification, \citet[Theorems~1.1 and~1.2]{BellExtremizers2026}
prove that every three-point maximizing representation satisfies $a=0$,
$\nu=0$, and $\operatorname{diam}(\supp X_i)\to0$; if infinitely many $X_i$
are nonconstant, then $S$ is nonatomic, and under $\beta\le18481/25000$ the
variances satisfy
\[
 \sum_{i\ge1}\Var(X_i)\operatorname{diam}(\supp X_i)^{-\gamma}<\infty
\]
for all $0<\gamma<\sqrt5-1$ (and all $0<\gamma<\sqrt{21}-3$ when all but
finitely many summands are two-point), ruling out two-sided $C^{1,1}$
regularity of $F_S$ at a maximizing threshold.

\subsection{Proved classifications of maximizers and the binomial phase law}\label{subsec:proved-classifications}

In Manuscripts~5 and~6 \citep{BellBernoulli2026,BellSmallStructure2026}, the
structural classification question (Question~\ref{q:structure}) is solved
completely in three regimes, and the asymptotic behavior of the optimal zero
padding in $P_n=\max_{1\le m\le n}B_m$ is determined:

\begin{enumerate}[label=(\arabic*),leftmargin=2em]
\item \emph{Common-diameter two-point arrays at every horizon $n\ge1$}
\citep[Theorem~1.1]{BellBernoulli2026}: Let $C_n^{\mathrm{com}}$ denote the
maximum of $\mathcal C(X_1,\ldots,X_n)$ over $n$-summand two-point arrays
whose nondegenerate members share a common support diameter $h>0$ (permitting
zero summands). Then for every $n\ge1$,
\begin{equation}\label{eq:cn-com-theorem}
 C_n^{\mathrm{com}}=P_n=\max_{1\le m\le n}B_m,
\end{equation}
and every common-diameter array attaining $C_n^{\mathrm{com}}$ consists of
$m\in\{1,\ldots,n\}$ identically distributed nondegenerate two-point summands
and $n-m$ zero summands. The proof expresses the convolution probabilities as
affine functions of the elementary symmetric polynomials of the success
probabilities $(p_1,\ldots,p_m)$ and combines Hoeffding's fiber theorem
\citep{Hoeffding1956} with certified negative-curvature enclosures along
constant-mean probability fibers.

\item \emph{Two two-point summands with arbitrary diameters ($n=2$)}
\citep[Theorem~1.2]{BellBernoulli2026}: Over arrays of at most two independent
two-point summands with arbitrary support diameters $h_1,h_2>0$, the maximum
of $\mathcal C(X_1,X_2)$ equals $C_1=B_1\approx0.3703515235$, and every
attaining array has exactly one nonconstant summand (so $C_2^{\mathrm{bin}}=P_2=B_1$).
If both summands are required to have strictly positive variance, the supremum
remains $C_1$ and is not attained.

\item \emph{Strict local improvement under common-diameter replacement}
\citep[Theorem~1.3]{BellBernoulli2026}: Fix $c_0>0$. There exist constants
$A_0,\sigma_0,\eta_0,\kappa>0$, depending only on $c_0$, with the following
property. Let $X_i=g_i(B_i-p_i)$ ($1\le i\le n$) be independent nondegenerate
two-point summands with $g_i>0$ and $0<p_i<1$, and set
\[
 a_i=p_i(1-p_i),\qquad
 A=\sum_{i=1}^n a_i,\qquad
 h_0=\frac1A\sum_{i=1}^n a_i g_i,\qquad
 e_i=\frac{g_i}{h_0}-1,
\]
\[
 \sigma=\left(\sum_{i=1}^n a_i e_i^2\right)^{1/2},\qquad
 \eta=\max_{1\le i\le n}|e_i|.
\]
If $A\ge A_0$, $0<\sigma\le\sigma_0$, $\eta\le\eta_0$, and
$\mathcal C(X_1,\ldots,X_n)\ge c_0$, then replacing all diameters $g_i$ by
their variance-weighted mean $h_0$ while keeping $(p_1,\ldots,p_n)$ fixed
strictly increases the Berry--Esseen ratio:
\begin{equation}\label{eq:local-diameter-gain}
 \mathcal C(B_1-p_1,\ldots,B_n-p_n)-\mathcal C(X_1,\ldots,X_n)
 \ge\kappa\frac{\sigma^2}{\sigma+\eta}>0.
\end{equation}
If in addition all $p_i$ lie in a fixed compact subset $J\subset(0,1)$, the
gain is at least $\kappa_J\sigma$, whereas when rare probabilities $p_i\to0$
are permitted the rate $\sigma^2/(\sigma+\eta)$ is sharp. (Note that the
strict positivity hypothesis $0<\sigma$ is essential, since $\sigma=0$ means
$g_1=\cdots=g_n=h_0$ where the two sides of \eqref{eq:local-diameter-gain}
are equal.)

\item \emph{Asymptotic lattice structure and exact classification for
small Lyapunov ratio} \citep[Theorems~1.1 and~1.2]{BellSmallStructure2026}:
First, along any sequence of finite independent arrays with $L_n\to0$ and
$\mathcal C_n\to\ce$, there exist lattice spacings
$h_n\sim R_n/\bigl((\pe^2+\qe^2)V_n\bigr)$ and centered laws $X_{n,i}^{\mathrm{lat}}$
supported on cosets of $h_n\Z$ such that
\[
 \sum_i W_2^2(X_{n,i},X_{n,i}^{\mathrm{lat}})=o(h_n^2),\qquad
 \sum_i W_3^3(X_{n,i},X_{n,i}^{\mathrm{lat}})=o(h_n^3),
\]
while $\mathcal B(X_n)\to1$ and a core of
$m_n\sim(\pe^2+\qe^2)^2/(\pe\qe L_n^2)$ summands carries $1-o(1)$ of $V_n$ and
$R_n$ and is $W_3$-close to scaled Esseen two-point laws. Second, there exists
an absolute constant $\eps_0>0$
(established by compactness and not effective) such that \emph{every}
independent array $(X_1,\ldots,X_n)$ attaining the unrestricted $n$-summand
maximum $C_n$ with $L<\eps_0$ consists of $m\in\{1,\ldots,n\}$ identically
distributed nondegenerate two-point summands and $n-m$ zero summands. To
eliminate arbitrarily small rounding residuals $T=\sum_i Z_i$ modulo $h_n\Z$
that cannot be ruled out by asymptotic $o(h_n^2)$ bounds alone, the proof
establishes a relative conditional first-moment bound
$\E(|T|\mid K=k)\ge c\,\E|T|$ valid for arbitrarily small $\E|T|>0$, proving
that exact lattice rounding strictly increases $\mathcal C$ whenever $T\ne0$;
discrete Green-function concavity on $h_n\Z$ then reduces each summand to two
adjacent lattice points, where \eqref{eq:cn-com-theorem} completes the
classification.

\item \emph{Asymptotic phase expansion of $B_m$ and optimal zero summands}
\citep[Theorem~4.1]{BellBernoulli2026}: A uniform seventh-order
Edgeworth--Euler--Maclaurin expansion of $R_m(p)$ around $p=\pe$ shows that
as $m\to\infty$,
\begin{equation}\label{eq:bm-phase-expansion}
 B_m=\ce-\frac{Q_*}{m}
 -\frac{H_*+\alpha\dist(m\pe-\delta_0,\Z)^2}{m^2}+O(m^{-3}),
\end{equation}
where $H_*$ is an explicit constant independent of $m$, and
\begin{equation}\label{eq:phase-constants}
 Q_*=\frac{\phz(867+287\sqrt{10})}{25920},\qquad
 \alpha=\frac{\phz(60+19\sqrt{10})}{30},\qquad
 \delta_0=\frac{2590+577\sqrt{10}}{12960}.
\end{equation}
Consequently, when comparing $m=n-j$ active summands ($j\ge0$ bounded)
against $n$, choosing $j$ zero summands minimizes the second-order loss
\begin{equation}\label{eq:parabola-envelope}
 \eta j+\dist(\theta_n+j\pe,\Z)^2,\qquad
 \eta=\frac{Q_*}{\alpha}=\frac{2510-747\sqrt{10}}{8640}\approx0.017104003,
\end{equation}
where $\theta_n=\operatorname{round}(n\pe-\delta_0)-(n\pe-\delta_0)\in[-1/2,1/2]$.
Because $5\eta$ strictly exceeds the maximum over $\theta\in[-1/2,1/2]$ of the
lower envelope of the first five parabolas $j\in\{0,1,2,3,4\}$, at most four
zero summands are ever optimal for sufficiently large $n$. Extending the
expansion to $O(m^{-3})$ and using a quadratic-field norm bound in
$\mathbb Q(\sqrt{10})$ to rule out exact ties between distinct branches
$j\ne j'$ proves that for all sufficiently large $n$:
\begin{enumerate}[label=(\alph*),leftmargin=2em]
\item the maximizing active count $m_n$ in $P_n=\max_{1\le m\le n}B_m$ is
unique;
\item the number of zero summands $j_n=n-m_n$ belongs to $\{0,1,2,3,4\}$, and
by Weyl equidistribution of the irrational rotation $n\pe\bmod1$ the
asymptotic frequencies of $j_n=0,1,2,3,4$ are respectively
\begin{equation}\label{eq:padding-frequencies}
 \pe-\omega,\quad
 \pe-\omega,\quad
 \sqrt{10}-3,\quad
 \omega,\quad
 \omega,
\end{equation}
where $\omega=(7773-2447\sqrt{10})/8640\approx0.00404464$;
\item the gaps between successive record active counts $m$ (those $m$ for
which $B_m>P_{m-1}$) belong to $\{2,3,5\}$.
\end{enumerate}
\end{enumerate}

\section{Computational evidence and certified local obstructions}\label{sec:computations}

\subsection{Scope of the stored computations}\label{subsec:comp-scope}

To interpret the numerical evidence in this repository accurately, one must
distinguish three different kinds of computation:
\begin{enumerate}[label=(\arabic*),leftmargin=2em]
\item \emph{High-precision floating-point evaluations in low-dimensional
parametric families:} The file \texttt{data/\allowbreak empirical-observations.json}
records deterministic evaluations of the one-variable binomial ratio $R_m(p)$
for $1\le m\le3000$ and of finite binomial-plus-compound-Poisson families.
These computations evaluate explicit closed-form or series expressions at
high precision within specified parametric families; they are neither interval
root-count certificates nor searches over arbitrary $n$-summand distributions.
\item \emph{Rigorous interval arithmetic certificates:} The certificate file
\texttt{certificates/\allowbreak observations/\allowbreak local-trap.json},
verified by \texttt{verify/\allowbreak observations/\allowbreak local\_trap.py}
in 320-bit Arb ball arithmetic \citep{Johansson2017,Johansson2018} (alongside
the interval certificates for Manuscripts~1, 2, and~5 under
\texttt{certificates/}), provides outward-rounded
computer-assisted mathematical proofs of specific inequalities, root
enclosures, and Hessian sign conditions.
\item \emph{Unrestricted multi-summand numerical optimization:} No general
numerical search over unrestricted multi-atom or arbitrary-diameter
$n$-summand arrays is included in the verified repository records, and no
claim in this paper or the companion manuscripts relies on uncertified
exploratory searches.
\end{enumerate}

\subsection{High-precision binomial evaluations and zero padding through $m=3000$}\label{subsec:binomial-data}

The dataset \texttt{data/\allowbreak empirical-observations.json} evaluates the i.i.d.\
binomial maxima $B_m=\max_{0<p<1}R_m(p)$ and the padded binomial envelope
$P_n=\max_{1\le m\le n}B_m$ for $1\le m\le3000$ in two regimes:
\begin{itemize}[leftmargin=2em]
\item \emph{Exhaustive threshold evaluation ($1\le m\le100$):} For each
$m\in\{1,\ldots,100\}$, a uniform grid of $4096$ probabilities in $(0,1/2]$
(which covers $(0,1)$ by reflection $p\leftrightarrow 1-p$) evaluates every
threshold index $0\le k\le m$ and both one-sided CDF branches in
\eqref{eq:rmp}. Visible local maxima are bracketed and refined by golden-section
and stationary-root iteration to $60$ decimal places, with stationarity
residuals at most $2.61\times10^{-59}$.
\item \emph{Central-threshold scan ($101\le m\le3000$):} For each
$m\in\{101,\ldots,3000\}$, the seven central threshold indices $k$ nearest
$mp$ are evaluated and bracketed at cell endpoints. As a consistency check,
for $101\le m\le450$ a denser control scan evaluates all central threshold
cells, both CDF branches, and $17$-point sub-meshes per cell with zero
discrepancies; for ratios exceeding $0.4045$, the tail bounds of
\citet{Schulz2016} rigorously exclude noncentral thresholds.
\end{itemize}

\begin{table}[t]
\centering
\caption{Selected values of the padded binomial envelope
$P_n=\max_{1\le m\le n}B_m=B_{m_n}$ from
\texttt{data/empirical-observations.json},
showing the maximizing active count $m_n$, the number of zero summands
$j_n=n-m_n$, the optimal success probability $p<\frac12$, the normalized
support diameter $h=1/\sqrt{m_n p(1-p)}$, the Lyapunov ratio $L$, and the
ratio $P_n$. Values for $m_n\le100$ and $m_n\in\{409,412,413\}$ are rounded
from 60-digit stationary-root refinements; $n=414$ and $n=3000$ are binary64
central-cell evaluations.}
\label{tab:binomial-optima}
\smallskip
\small
\begin{tabular}{@{}rrrllll@{}}
\toprule
$n$ & $m_n$ & $j_n=n-m_n$ & Probability $p$ & Diameter $h$ & Lyapunov $L$ & Ratio $P_n=B_{m_n}$ \\
\midrule
1    & 1    & 0 & 0.390958631281333 & 2.049327 & 1.073396 & 0.370351523474379 \\
2    & 1    & 1 & 0.390958631281333 & 2.049327 & 1.073396 & 0.370351523474379 \\
3    & 3    & 0 & 0.431607387750433 & 1.165657 & 0.593733 & 0.398156326353304 \\
4    & 3    & 1 & 0.431607387750433 & 1.165657 & 0.593733 & 0.398156326353304 \\
10   & 10   & 0 & 0.429865759681911 & 0.638771 & 0.325669 & 0.406507736905768 \\
12   & 10   & 2 & 0.429865759681911 & 0.638771 & 0.325669 & 0.406507736905768 \\
50   & 49   & 1 & 0.414895409236786 & 0.289945 & 0.149173 & 0.409145180658100 \\
100  & 99   & 1 & 0.417421147624576 & 0.203806 & 0.104683 & 0.409451697372009 \\
409  & 409  & 0 & 0.418879394858355 & 0.100221 & 0.051430 & 0.409665276377194 \\
412  & 409  & 3 & 0.418879394858355 & 0.100221 & 0.051430 & 0.409665276377194 \\
413  & 409  & 4 & 0.418879394858355 & 0.100221 & 0.051430 & 0.409665276377194 \\
414  & 414  & 0 & 0.418652506346914 & 0.099622 & 0.051129 & 0.409666049802979 \\
3000 & 2999 & 1 & 0.418913502013957 & 0.037011 & 0.018992 & 0.409723068705426 \\
\bottomrule
\end{tabular}
\end{table}

Table~\ref{tab:binomial-optima} displays representative values of $P_n$ and
the optimal zero padding $j_n=n-m_n$. Several features merit emphasis:
\begin{enumerate}[label=(\arabic*),leftmargin=2em]
\item Across the entire scan $1\le n\le3000$, the observed number of zero
summands $j_n=n-m_n$ takes only the five values $\{0,1,2,3,4\}$ proved
asymptotically in \citet[Theorem~4.1]{BellBernoulli2026}. At small horizons,
successive record active counts $m$ (those $m$ with $B_m>P_{m-1}$) include
consecutive integers with gap $1$ (such as $m=5,6$, $m=17,18$, and $m=29,30$)
before settling into the asymptotic gap set $\{2,3,5\}$ proved for large $n$
in \citet[Theorem~4.1]{BellBernoulli2026}.
\item The horizons $n=412$ and $n=413$ are the first values of $n$ at which
three and four zero summands occur: both retain the active count $m=409$,
whose 60-digit refined ratio compares with $B_{412}$ and $B_{413}$ as
\begin{align*}
 B_{409}&=0.40966527637719430986\ldots,\\
 B_{412}&=0.40966527592709729990\ldots\quad(B_{409}-B_{412}\approx 4.501\times10^{-10}),\\
 B_{413}&=0.40966478482671455982\ldots\quad(B_{409}-B_{413}\approx 4.916\times10^{-7}).
\end{align*}
\item Comparing the empirical padding $j_n$ with the leading quadratic phase
rule \eqref{eq:parabola-envelope} illustrates the role of the $O(m^{-3})$
error term in \eqref{eq:bm-phase-expansion} near the narrow phase windows of
width $\omega\approx0.004045$ where $j\in\{3,4\}$. Across the transition
block $n=409,410,411,412,413,414$ (where $j_n=0,1,2,3,4,0$, respectively) and
at $n=2957,2958$ (where $j_n=3,4$), the leading phase formula
\eqref{eq:parabola-envelope} already agrees with the observed padding $j_n$.
At smaller or borderline horizons such as $n\in\{147,148\}$, $n\in\{677,678\}$,
and $n\in\{1780,1781\}$, the $O(1)$ leading phase formula lies near the
boundary of the $\{3,4\}$ window, and including the finite-$m$ probability
drift shifts the maximizer to $j_n\in\{1,2\}$ (for $n=147,148$ and
$1780,1781$) or $j_n\in\{0,1\}$ (for $n=677,678$), matching the observed
active counts.
\end{enumerate}

\subsection{A certified strict local maximizer with unequal diameters}\label{subsec:local-trap}

Why does the local diameter-equalization theorem
(Theorem~1.3 of \citealp{BellBernoulli2026}, recalled in
\eqref{eq:local-diameter-gain}) require the diameter spread $\sigma$ and
maximum relative deviation $\eta$ to be small? When all nondegenerate
two-point summands share a common diameter $h$, the $2^m$ atoms of the
convolution sum $S=\sum_{i=1}^m h(B_i-p_i)$ coalesce into $m+1$ lattice points.
Perturbing $h_i$ splits the colliding atoms at the maximizing threshold and
produces a downward directional cusp in $\mathcal C$. However, atom collisions
also occur when unequal diameters have \emph{rational ratios}, and such
collisions can lock unequal diameters into a strict local maximum.

Consider $n=4$ independent two-point summands $X_i=h_i(B_i-p_i)$ with support
diameters
\begin{equation}\label{eq:trap-diameters}
 h=(1,1,1,2),
\end{equation}
upper atom probabilities $(p_1,p_2,p_3,p_4)=(p,p,p,r)$ near
\begin{equation}\label{eq:trap-probs}
 p=0.572751600606024506947692808\ldots,\qquad
 r=0.337472966089985017834868704\ldots,
\end{equation}
and raw threshold $k=2$ on the closed branch
$\Pp(\sum_{i=1}^4 h_i B_i\le 2)-\Phi(z_2)$, where
$z_2=(2-3p-2r)/\sqrt{3p(1-p)+4r(1-r)}$. At the raw sum value
$\sum_{i=1}^4 h_i B_i=2$, four binary configurations $v=(B_1,B_2,B_3,B_4)\in\{0,1\}^4$
collide:
\begin{equation}\label{eq:collision-vertices}
 v^{(1)}=(1,1,0,0),\qquad
 v^{(2)}=(1,0,1,0),\qquad
 v^{(3)}=(0,1,1,0),\qquad
 v^{(4)}=(0,0,0,1).
\end{equation}
These four vectors satisfy $v^{(j)}\cdot h=2$ and form the vertices of a
nondegenerate three-dimensional tetrahedron in the hyperplane
$\{v\in\R^4:v\cdot h=2\}$.

Using 320-bit Arb ball arithmetic on the closed box of radius $10^{-28}$ around
\eqref{eq:trap-probs}, the verification script
\texttt{verify/\allowbreak observations/\allowbreak local\_\allowbreak trap.py}
checks the interval certificate
\texttt{certificates/\allowbreak observations/\allowbreak local-\allowbreak trap.json}
and establishes the following theorem:

\begin{theorem}[Certified unequal-diameter strict local maximizer]\label{thm:local-trap}
In the four-summand two-point family $X_i=h_i(B_i-p_i)$ with $h=(1,1,1,2)$,
the closed box of radius $10^{-28}$ around $(p,r)$ in \eqref{eq:trap-probs}
contains a unique point $(p_*,r_*)$ at which the probabilities
$p^{(0)}=(p_*,p_*,p_*,r_*)$ are stationary for the closed-branch ratio at raw
threshold $k=2$. At $(h;p^{(0)})$:
\begin{enumerate}[label=(\roman*),leftmargin=2em]
\item The Kolmogorov distance $\Delta(S)$ is uniquely attained at the raw
threshold $k=2$ on the closed branch, and the Berry--Esseen ratio satisfies
\begin{equation}\label{eq:trap-ratio}
 \mathcal C\in 0.28248894967964153979405802\pm 5.54\times10^{-27}.
\end{equation}
\item The $4\times4$ Hessian $H_p=\nabla_p^2\mathcal C(h;p^{(0)})$ with
respect to the probabilities $(p_1,p_2,p_3,p_4)$ at fixed diameters
$h=(1,1,1,2)$ is strictly negative definite.
\item Under any perturbation $h\mapsto h+d$ ($d\in\R^4$ small) of the support
diameters, the four colliding configurations \eqref{eq:collision-vertices} move
to the raw sum values $v^{(j)}\cdot(h+d)=2+v^{(j)}\cdot d$ ($1\le j\le 4$). Any
candidate threshold near $2$ that retains all four colliding configurations in
the closed ray has the form $2+v^{(j_*)}\cdot d$ for an active index
$j_*\in\{1,2,3,4\}$ maximizing $v^{(j)}\cdot d$ (any lower threshold omits at
least one colliding configuration of positive probability). Along each smooth
branch function $f_j(h+d,p^{(0)}+\delta p)$ corresponding to the threshold
$v^{(j)}\cdot(h+d)=2+v^{(j)}\cdot d$, the diameter gradient at $(0,p^{(0)})$ is
$c_1(\beta-v^{(j)})$ with $c_1>0$, where $\beta=\sum_{j=1}^4 w_j v^{(j)}$ lies
in the strict interior of the tetrahedron
$\operatorname{conv}\{v^{(1)},v^{(2)},v^{(3)},v^{(4)}\}$ with barycentric
weights $w_1=w_2=w_3\in 0.31153\pm 10^{-5}$ and $w_4\in 0.06541\pm 10^{-5}$.
Consequently, for every $d=d_\perp$ orthogonal to $h=(1,1,1,2)$ and any active
index $j_*$ maximizing $v^{(j)}\cdot d_\perp$,
\begin{equation}\label{eq:trap-cusp}
 \nabla_d f_{j_*}(0,p^{(0)})\cdot d_\perp
 =c_1\Bigl(\beta\cdot d_\perp-\max_{1\le j\le 4}v^{(j)}\cdot d_\perp\Bigr)
 \le -c_2\|d_\perp\|,
\end{equation}
with $c_2>0$. Since $f_j(0,p)=\mathcal C(h;p)$ has vanishing probability
gradient $\nabla_p\mathcal C(h;p^{(0)})=0$ and negative-definite probability
Hessian $H_p\preceq -c_H I<0$, a second-order Taylor expansion of $f_{j_*}$
around $(0,p^{(0)})$ with $a=\|d_\perp\|$ and $b=\|\delta p\|$ gives
\[
 \mathcal C(h+d_\perp;p^{(0)}+\delta p)-\mathcal C(h;p^{(0)})
 \le -c_2 a-\frac{c_H}{2}b^2+O(ab+a^2+b^3)<0
\]
for all sufficiently small $(d_\perp,\delta p)\ne(0,0)$, because
$ab\le b(a+b^2)=o(a+b^2)$ and $a^2+b^3=o(a+b^2)$ as $a,b\to0$.
\item For both the unit-diameter summands ($i=1,2,3$) and the diameter-$2$
summand ($i=4$), the single-coordinate variational contact cubic majorants
hold globally on $\R$, with certified strictly positive gaps at all
non-contact breakpoints of the deleted-sum distribution functions.
\end{enumerate}
\end{theorem}

Theorem~\ref{thm:local-trap} shows that $(h;p^{(0)})$ is a strict local
maximizer of $\mathcal C$ under joint perturbations of $(h_1,h_2,h_3,h_4)$ and
$(p_1,p_2,p_3,p_4)$ modulo common scaling, and satisfies all single-coordinate
variational contact inequalities on $\R$. Yet its Berry--Esseen ratio
$0.282489\ldots$ is substantially smaller than the ratio
$\sqrt2/4\approx0.353553$ of two symmetric Bernoulli summands and two zeros
(evaluated at threshold $0$), and far below the four-summand common-diameter
maximum $P_4=B_3\approx0.398156$. This example demonstrates that local
stationarity and single-coordinate contact conditions cannot by themselves
force support diameters to be equal when $n\ge3$.

\subsection{Binomial sums with a compound-Poisson jump component}\label{subsec:cp-data}

The dataset \texttt{data/\allowbreak empirical-observations.json} (under
\texttt{mixed\_core\_cp}) also records numerical comparisons for centered sums
of the form
\begin{equation}\label{eq:mixed-cp-law}
 S_{\mathrm{raw}}=\bigl(\Bin(M,p)-Mp\bigr)+g\bigl((N_+-\ell_+)-(N_--\ell_-)\bigr),
\end{equation}
where $N_\pm\sim\Pois(\ell_\pm)$ are independent of $\Bin(M,p)$, across $18$
parameter profiles with $M\in\{1,2,3,5,10,20\}$, jump step $g\in\{1/2,1,2\}$,
$p\in[0.05,0.5]$, and total intensity $10^{-6}\le\ell_++\ell_-\le12$. In this
experiment, the variance and the third-moment functional are taken to be
\begin{equation}\label{eq:mixed-cp-price}
 V=Mp(1-p)+g^2(\ell_++\ell_-),\qquad
 \beta=Mp(1-p)\bigl(p^2+(1-p)^2\bigr)+g^3(\ell_++\ell_-).
\end{equation}
As emphasized in Remark~\ref{rem:balanced-and-id} and
Section~\ref{subsec:variational}, the term $g^3(\ell_++\ell_-)$ in
\eqref{eq:mixed-cp-price} is the L\'evy jump integral
$\int_{\R\setminus\{0\}}|y|^3\,\nu(dy)$ representing the limiting sum of third
absolute moments of arbitrarily many rare two-point summands of step sizes
$\pm g$ in the all-horizon compactification \eqref{eq:compactified-rep}. It is
not the third absolute moment of a single centered Poisson variable, and by
Remark~\ref{rem:balanced-and-id}(2) it is strictly smaller than the minimum
third-moment sum $\mathcal M_{n-M}$ of any finite $(n-M)$-summand independent
factorization of the compound-Poisson component.

Even when the compound-Poisson component is charged only this smaller limiting
third-moment cost $g^3(\ell_++\ell_-)$, across all $18$ tested profiles every
candidate with $\ell_++\ell_->0$ has ratio $\Delta(S_{\mathrm{raw}})/(\beta/V^{3/2})$
strictly smaller than the pure binomial ratio $R_M(p)$ ($\ell_+=\ell_-=0$) at
the same probability $p$. This provides a finite-$M$ numerical illustration of
the theorem $\nu=0$ proved for compactified global maximizers in
\citet[Theorem~1.1]{BellExtremizers2026}.

\section{What remains open}\label{sec:open}

\subsection{The finite-$n$ conjectures}\label{subsec:open-conjectures}

Combining the theoretical results of Sections~\ref{sec:factorization}--\ref{sec:known}
and the numerical evidence of Section~\ref{sec:computations}, we can state the
three open finite-$n$ conjectures corresponding to
Questions~\ref{q:value}--\ref{q:uniqueness}.

\begin{conjecture}[Finite-$n$ value conjecture]\label{conj:value}
For every integer $n\ge1$,
\[
 C_n=P_n=\max_{1\le m\le n}B_m.
\]
In particular, $\Cind=\sup_{n\ge1}C_n=\lim_{m\to\infty}B_m=\ce$.
\end{conjecture}

\begin{conjecture}[Finite-$n$ extremizer classification conjecture]\label{conj:structure}
For every integer $n\ge1$, every $n$-summand independent array
$(X_1,\ldots,X_n)$ attaining $C_n$ consists of $m\in\{1,\ldots,n\}$
identically distributed nondegenerate two-point random variables and $n-m$
zero variables, up to a common nonzero scale factor, sign reflection, and
permutation of coordinates.
\end{conjecture}

\begin{conjecture}[Uniqueness of the binomial maximizer]\label{conj:uniqueness}
\begin{enumerate}[label=(\roman*),leftmargin=2em]
\item For every integer $m\ge1$, the function $p\mapsto R_m(p)$ in
\eqref{eq:rmp} attains its maximum $B_m$ on $(0,1/2]$ at a unique probability
$p_m\in(0,1/2)$.
\item For every integer $n\ge1$, the padded binomial envelope
$P_n=\max_{1\le m\le n}B_m$ is attained at a unique active count
$m_n\in\{1,\ldots,n\}$.
\end{enumerate}
\end{conjecture}

\begin{table}[t]
\centering
\caption{Summary of what is currently proved toward Conjectures~\ref{conj:value}--\ref{conj:uniqueness}.}
\label{tab:status-summary}
\smallskip
\small
\begin{tabular}{@{}lll@{}}
\toprule
Regime & Proved result & Reference \\
\midrule
$n=1$, unrestricted & $C_1=B_1$; unique two-point maximizer & \citet{BentkusKirsa1989} \\
$n=2$, two-point & $C_{2,\mathrm{bin}}=P_2=B_1$; one summand zero & \citet[Thm.~1.2]{BellBernoulli2026} \\
$n=2$, unrestricted & $C_2<49/120<\ce$ ($C_2=B_1$ open) & \citet[Thm.~1.1]{BellFinite2026} \\
$n=3$, two-point & $\Delta\le\ce L$ ($\sup\mathcal C=B_3$ open) & \citet[Thm.~1.2(i)]{BellFinite2026} \\
All $n\ge1$, common-diam.\ two-point & $C_n^{\mathrm{com}}=P_n$; maximizers padded i.i.d. & \citet[Thm.~1.1]{BellBernoulli2026} \\
All $n\ge1$, $L\le\exp(-\exp(97))$ & $\Delta\le\ce L$ & \citet[Thm.~1.1]{BellSmall2026} \\
All $n\ge1$, maximizers with $L<\eps_0$ & $C_n=P_n$; maximizers padded i.i.d. & \citet[Thm.~1.2]{BellSmallStructure2026} \\
All $n\ge1$, unrestricted & $\Cind<0.44988$; $\le4$ pts ($\le3$-pt rep.) & \citet{BellIndependent2026,BellExtremizers2026} \\
Large $n$, binomial envelope $P_n$ & Unique $m_n$; $n-m_n\in\{0,1,2,3,4\}$ & \citet[Thm.~4.1]{BellBernoulli2026} \\
\bottomrule
\end{tabular}
\end{table}

Table~\ref{tab:status-summary} summarizes the proved cases. Note that
Conjecture~\ref{conj:structure} at a given $n$ immediately implies
Conjecture~\ref{conj:value} at that $n$.

\subsection{Overall proof architecture and the missing global implication}\label{subsec:missing-implication}

The results of this paper and the companion manuscripts assemble the
machinery for two complementary approaches to Conjecture~\ref{conj:structure},
and make it transparent which implications are established and which remain
open in each:

\begin{enumerate}[label=\textbf{Route \arabic*.},ref=\arabic*,leftmargin=5.5em]
\item\label{route:sum-law} \textbf{(Reduction to sum-law selection via exact-fiber factorization)}
In Section~\ref{sec:factorization}, we separated the optimization of
$\mathcal C(X_1,\ldots,X_n)$ into two stages:
\[
 C_n=\sup_{\mu}\mathcal Q_n(\mu),\qquad
 \mathcal Q_n(\mu)=\frac{V(\mu)^{3/2}\Delta(\mu)}{\mathcal M_n(\mu)},
\]
where $\mu=\Law(S)$ is the distribution of the sum and $\mathcal M_n(\mu)$ is
the minimum third-moment sum over all $n$-summand independent factorizations
of $\mu$. By Theorem~\ref{thm:binomial-fiber}\ref{it:binomial-equivalence},
once $\mu$ is known to be a centered binomial law $\Law(g(\Bin(m,p)-mp))$ of
order $m\le n$, minimizing $\sum_{i=1}^n\E|X_i|^3$ on the fiber
$\mathcal F_n(\mu)$ \emph{automatically} forces the summands $(X_1,\ldots,X_n)$
to consist of $m$ identically distributed two-point variables and $n-m$ zeros.
More generally, by Theorem~\ref{thm:binomial-fiber}\ref{it:pb-fiber}, if
$\mu$ is any Poisson-binomial law of order $k\le n$, fiber minimization forces
every nonconstant summand to be a two-point variable of common diameter $g$,
and Theorem~1.1 of \citet{BellBernoulli2026} then forces their probabilities
$p_j$ to be equal.

\emph{What remains unproved in Route~\ref{route:sum-law}:} Proving that the
sum distribution $\mu=\Law(S)$ of every global maximizer of $C_n$ is a
Poisson-binomial (or binomial) law of order at most $n$. Minimizing the
factorization cost $\mathcal M_n(\mu)$ for a fixed $\mu$ operates within a
single convolution fiber $\mathcal F_n(\mu)$; it does not compare
$\mathcal Q_n(\mu)$ across different sum distributions $\mu$.

\item\label{route:coherence} \textbf{(Reduction to characteristic-function coherence)}
In Section~\ref{sec:criteria}, Theorem~\ref{thm:fourier-criterion} showed that
for any independent array $(X_1,\ldots,X_n)$, the moment-determined span
$\widehat h=(R+D)/V$ defines a nonnegative characteristic-function deficit
$\mathcal J=1-|\E e^{2\pi\ii S/\widehat h}|^2\ge0$ such that $\mathcal J=0$ if
and only if $(X_1,\ldots,X_n)$ is a padded i.i.d.\ two-point array. Notice
that $\mathcal J=0$ combines two properties simultaneously: it places the
support of every summand on a coset of $\widehat h\Z$ and saturates the
separated-support moment inequality of Theorem~\ref{thm:separated-moment} at
the same span $\widehat h$. (By contrast, merely knowing that the summands of
a fixed-$n$ maximizer are supported on cosets of some common lattice $h\Z$
does not by itself reduce them to two-point laws without an additional step:
the discrete Green-function concavity argument of
\citealp[Section~9]{BellSmallStructure2026} requires $L<\eps_0$ to supply
asymptotic local limit estimates for the deleted sum, while
Theorem~\ref{thm:separated-moment} requires equality in the moment bound
\eqref{eq:array-moment-ineq} in addition to support separation $h$.)

\emph{What remains unproved in Route~\ref{route:coherence}:} Proving that
global optimality for $\mathcal C(X_1,\ldots,X_n)$ forces $\mathcal J=0$. As
emphasized in Remark~\ref{rem:j-not-monotone}, Theorem~\ref{thm:fourier-criterion}
is an equality characterization for $\mathcal J=0$, not an inequality bounding
$\mathcal C$ by a decreasing function of $\mathcal J$. Moreover, while
\citet[Theorem~3.2]{BellExtremizers2026} guarantees that every fixed-$n$
maximizer has at most four support points per summand and admits an equally
maximizing replacement with at most three support points per summand (which is
convolution-indecomposable when at least one zero slot is present), proving
the value conjecture $C_n=P_n$ by this route suffices if at least one
three-point maximizing representative for each $n$ satisfies $\mathcal J=0$,
whereas proving the full classification conjecture
(Conjecture~\ref{conj:structure}) requires showing that every three-point
maximizing representative satisfies $\mathcal J=0$ and that no four-point
flat-segment maximizer can exist.
\end{enumerate}

\subsection{Specific analytical obstacles in the remaining regimes}\label{subsec:obstacles}

Three concrete analytical phenomena explain why extending the proved
classifications to all $n\ge2$ (for general summands) and all $n\ge3$ (for
two-point summands with arbitrary diameters) requires new global arguments:

\begin{enumerate}[label=(\arabic*),leftmargin=2em]
\item \emph{Zero summands decouple the horizon $n$ from the Lyapunov ratio $L$:}
Theorem~1.2 of \citet{BellSmallStructure2026} proves
Conjectures~\ref{conj:value} and~\ref{conj:structure} for any maximizer of
$C_n$ whose Lyapunov ratio satisfies $L<\eps_0$. In the i.i.d.\ setting,
$L_n=(p^2+q^2)/\sqrt{npq}\to0$ automatically as $n\to\infty$. In the
independent setting, however, adjoining zero summands increases $n$ without
changing $L$: an $n$-summand array can have $L\ge\eps_0$ for arbitrarily large
$n$ if all but a bounded number of summands have zero or negligible variance.
Because $\eps_0>0$ in \citet{BellSmallStructure2026} is obtained by a
compactness contradiction argument and is not effective, ruling out
maximizers with $L\ge\eps_0$ requires a finite-$n$ classification at moderate
Lyapunov ratios.

\item \emph{Atom collisions at rational diameter ratios create strict local
traps:}
For two-point summands with arbitrary support diameters $h=(h_1,\ldots,h_n)$,
Theorem~\ref{thm:local-trap} proves that at $n=4$ the unequal-diameter
configuration $h=(1,1,1,2)$ is a strict local maximizer of $\mathcal C$ that
satisfies every single-coordinate variational contact inequality globally on
$\R$. Whenever a subset of diameters has rational ratios, multiple binary
configurations can collide at the maximizing threshold and place the effective
weight vector $\beta$ in the interior of the collision polytope in
\eqref{eq:trap-cusp}, locking the diameters against all local perturbations.
Thus any proof that $C_n^{\mathrm{bin}}=P_n$ for $n\ge3$ must use a global
replacement (capable of jumping across collision hyperplanes or comparing the
multi-collision value against $P_n$) rather than local gradient ascent.

\item \emph{Dependence of single-summand replacements on the deleted-sum
distribution:}
When one modifies a single summand $X_i$ while holding the remaining summands
fixed, the numerator of $\mathcal C$ at a threshold $x$ is
$\bigl|\E F_{T_i}(x-X_i)-\Phi(x/\sqrt V)\bigr|$, where $F_{T_i}$ is the
distribution function of the deleted sum $T_i=\sum_{j\ne i}X_j$. Unlike the
third-moment sum $R=\sum_{i=1}^n\E|X_i|^3$, which is additive across
coordinates, $\E F_{T_i}(x-X_i)$ couples $X_i$ to the step locations and jump
sizes of $T_i$. As Proposition~\ref{prop:transport-counterexample}
illustrates for zero-bias transport, and as Theorem~\ref{thm:local-trap}
illustrates for unequal-diameter two-point arrays, a marginal replacement that
does not account for the alignment of atoms with the threshold $x$ and the
deleted sum $T_i$ can decrease $\mathcal C$ even when it moves $X_i$ toward a
two-point or equal-diameter law. Establishing an admissible global comparison
that overcomes this coupling is the central open problem in the finite-$n$
theory of the independent Berry--Esseen inequality.
\end{enumerate}

\bibliographystyle{plainnat}
\bibliography{refs}
\end{document}
